% =============================================================================
% Розділ 2. Програмне забезпечення та програмна реалізація системи
% Підключення: % =============================================================================
% Розділ 2. Програмне забезпечення та програмна реалізація системи
% Підключення: % =============================================================================
% Розділ 2. Програмне забезпечення та програмна реалізація системи
% Підключення: % =============================================================================
% Розділ 2. Програмне забезпечення та програмна реалізація системи
% Підключення: \input{sections/section2_implementation}
% Потрібно в преамбулі: \usetikzlibrary{arrows.meta,positioning,babel}
% =============================================================================
\providecommand{\px}{пкс}               % пікселі карти
\providecommand{\code}[1]{\texttt{#1}}  % імена файлів, модулів і функцій
\newpage
\section{Програмне забезпечення та програмна реалізація системи}
\label{sec:implementation}

\subsection{Постановка задачі та загальна архітектура системи}

Розроблена система розв'язує задачу абсолютної візуальної локалізації БпЛА: за кадром із бортової камери та ортофотопланом місцевості визначити, над якою точкою карти перебуває апарат, а за послідовністю кадрів --- відновити траєкторію польоту. Інформація про те, якому фрагменту карти відповідає кадр, системі не надається: відповідну ділянку ортофотоплану система знаходить автоматично.

Вхідними даними є:
\begin{itemize}
    \item ортофотоплан місцевості розміром $W \times H$ пікселів;
    \item відеозапис польоту БпЛА або послідовність кадрів із часовими мітками.
\end{itemize}

Вихідними даними є:
\begin{itemize}
    \item для кожного кадру --- координати $(X, Y)$ проєкції центра кадру на ортофотоплан і контур огляду камери (чотирикутник на карті);
    \item для послідовності кадрів --- траєкторія польоту з позначкою статусу кожного кадру та показниками якості зіставлення.
\end{itemize}

Робота системи ґрунтується на таких припущеннях:
\begin{itemize}
    \item камера спрямована приблизно вертикально вниз (надирна зйомка);
    \item сцену в межах кадру можна наближено вважати плоскою, тому перетворення <<кадр $\to$ карта>> описується гомографією \cite{hartley};
    \item відношення просторової роздільності кадру та карти відоме з точністю до множника $\approx 2$ (див. п.~\ref{sec:scale});
    \item ортофотоплан не має географічної прив'язки, тому всі координати та похибки обчислюються в пікселях карти (\px{}), без переведення в метри.
\end{itemize}

Загальну схему обробки наведено на рисунку~\ref{fig:pipeline}. Система складається з офлайн-етапу підготовки карти, який виконується один раз, та онлайн-етапу обробки кадрів. Ключовим елементом онлайн-етапу є єдиний конвеєр зіставлення пари <<кадр --- фрагмент карти>>, який однаково застосовується в усіх режимах роботи: під час перевірки відомої пари, автоматичного пошуку фрагмента та трекінгу послідовності.

\begin{figure}[H]
\centering
\begin{tikzpicture}[
    node distance=5mm and 6mm,
    box/.style={draw, rounded corners=2pt, align=center, font=\small,
                minimum width=42mm, minimum height=9mm, inner sep=3pt},
    off/.style={box, fill=gray!12},
    on/.style={box, fill=blue!7},
    dec/.style={box, fill=orange!12},
    arr/.style={-{Stealth[length=2.2mm]}, thick}]
  % офлайн-етап
  \node[off] (map) {Ортофотоплан\\$6200\times4320$};
  \node[off, below=of map] (tiles) {Розбиття на фрагменти\\$512\times512$, перекриття 50\,\%};
  \node[off, below=of tiles] (cache) {Кеш ознак фрагментів\\(SuperPoint)};
  % онлайн-етап
  \node[on, right=16mm of map] (frame) {Кадр БпЛА\\$1920\times1080$};
  \node[on, below=of frame] (prep) {Нормалізація масштабу\\і компенсація курсу};
  \node[dec, below=of prep] (pool) {Пул фрагментів:\\вся карта / вікно $3\times3$};
  \node[on, below=of pool] (match) {Зіставлення\\LightGlue / LoFTR};
  \node[on, below=of match] (geom) {RANSAC + гомографія,\\перевірка правдоподібності};
  \node[dec, below=of geom] (rank) {Фільтр і ранжування\\фрагментів};
  \node[on, below=of rank] (coord) {Координати $(X,Y)$\\і контур огляду};
  \node[dec, below=of coord] (track) {Кінематична перевірка,\\траєкторія};
  \draw[arr] (map) -- (tiles);
  \draw[arr] (tiles) -- (cache);
  \draw[arr] (frame) -- (prep);
  \draw[arr] (prep) -- (pool);
  \draw[arr] (pool) -- (match);
  \draw[arr] (match) -- (geom);
  \draw[arr] (geom) -- (rank);
  \draw[arr] (rank) -- (coord);
  \draw[arr] (coord) -- (track);
  \draw[arr] (cache.south) |- (match.west);
  \draw[arr, dashed] (track.east) -- ++(6mm,0) |- node[pos=0.25, right, font=\footnotesize, align=left]{прогноз\\положення\\і курсу} (pool.east);
\end{tikzpicture}
\caption{Загальна схема роботи системи візуальної локалізації (сірим позначено офлайн-етап)}
\label{fig:pipeline}
\end{figure}

\subsection{Використані програмні засоби}

Систему реалізовано мовою Python~3.11. Використано такі бібліотеки:
\begin{itemize}[label=\textbf{•}]
    \item \textbf{PyTorch} 2.14 з підтримкою CUDA~12.6 \cite{pytorch} --- виконання нейромережевих моделей на GPU;
    \item \textbf{LightGlue} (офіційна реалізація авторів) --- моделі SuperPoint \cite{superpoint} і LightGlue \cite{lightglue};
    \item \textbf{Kornia} 0.8 \cite{kornia} --- реалізація LoFTR \cite{loftr} з попередньо навченими вагами \textit{outdoor};
    \item \textbf{OpenCV} 5.0 \cite{opencv} --- читання та масштабування зображень, оцінювання гомографії методами RANSAC \cite{ransac} і MAGSAC++ \cite{magsac}, геометричні перетворення;
    \item \textbf{NumPy} і \textbf{pandas} --- числові обчислення, робота з таблицями метаданих і результатів;
    \item \textbf{Matplotlib} --- побудова графіків;
    \item \textbf{pytest} --- автоматичне тестування модулів.
\end{itemize}

Нейромережі в роботі не навчалися: використано опубліковані попередньо навчені ваги моделей.

\subsection{Тестове середовище}

Розроблення та експерименти виконувалися на ноутбуці з такими характеристиками:
\begin{itemize}[label=\textbf{--}]
    \item \textbf{процесор:} AMD Ryzen 5 4600H;
    \item \textbf{оперативна пам'ять:} 8 ГБ;
    \item \textbf{відеокарта:} NVIDIA GeForce GTX 1650 (4 ГБ відеопам'яті);
    \item \textbf{CUDA:} 12.6 (збірка PyTorch \texttt{cu126});
    \item \textbf{операційна система:} Windows 11.
\end{itemize}

\subsection{Вхідні дані}

Для моделювання використано ортофотоплан району розміром $6200 \times 4320$ пікселів (формат RGBA з повністю непрозорим альфа-каналом) і відеозапис польоту БпЛА DJI (файл \code{DJI\_0331.MP4}): роздільність $1920 \times 1080$, частота 29,97~кадр/с, тривалість $\approx 94$~с, кодек HEVC, камера спрямована вниз.

Відео та ортофотоплан знято в різні пори року: ортофотоплан --- влітку, відео --- взимку. Через це рослинність, тіні та покриття ґрунту на кадрах і на карті суттєво відрізняються, що ускладнює зіставлення.

\subsection{Підготовка ортофотоплану}

\textbf{Розбиття на фрагменти з перекриттям.}\\
Нейромережеві методи зіставлення розраховані на зображення розміром порядку сотень пікселів, тому ортофотоплан розбивається на квадратні фрагменти (тайли) розміром $S = 512$ пікселів із часткою перекриття $o = 0{,}5$. Крок сітки визначається як
\begin{equation}
s = \lfloor S \cdot (1 - o) \rfloor = 256 .
\end{equation}
Координати лівих верхніх кутів фрагментів уздовж осі $x$ утворюють послідовність $0, s, 2s, \dots$; якщо регулярна сітка не доходить до правого краю карти, додається ще один крайовий фрагмент із початком у $W - S$ (аналогічно для осі $y$). Завдяки цьому всі фрагменти мають однаковий розмір і кожен піксель карти покрито. Кількість фрагментів:
\begin{equation}
N = \left( \left\lceil \frac{W - S}{s} \right\rceil + 1 \right) \cdot \left( \left\lceil \frac{H - S}{s} \right\rceil + 1 \right) = 24 \cdot 16 = 384 .
\end{equation}
Для кожного фрагмента в таблицю метаданих \code{tiles.csv} записуються ідентифікатор і координати $(X_0, Y_0)$ лівого верхнього кута в пікселях повної карти. Ці координати є основою для переходу від локальних координат фрагмента до глобальних координат карти.

Перекриття 50\,\% гарантує, що будь-яка ділянка карти, менша за половину фрагмента, повністю потрапляє принаймні в один фрагмент. Водночас, як показано в розділі~\ref{sec:experiments}, контур огляду камери на карті має розмір $\approx 518 \times 292$~\px{} і є ширшим за фрагмент, тому кадр, як правило, лише частково перекривається з кожним окремим фрагментом.

\textbf{Перевірка цілісності даних.}\\
Після розбиття автоматично перевіряється, що кожен фрагмент попіксельно збігається з відповідною ділянкою вихідної карти, усі фрагменти лежать у межах карти, їх кількість відповідає формулі, а таблиця метаданих не містить пропусків і дублікатів.

\subsection{Вилучення кадрів із відео}

З відео вибирається один кадр на секунду: кожен $n$-й кадр, де $n = \operatorname{round}(f_{\text{відео}} / f_{\text{ціль}}) = \operatorname{round}(29{,}97 / 1) = 30$. Часова мітка кадру обчислюється як $t = i / f_{\text{відео}}$, де $i$ --- номер кадру у відео. Кадри декодуються послідовно (без позиціювання за номером, яке ненадійне для кодеків H.264/HEVC) і зберігаються без втрат у форматі PNG. У результаті отримано 95 кадрів із кроком $\approx 1{,}001$~с.

\subsection{Нормалізація масштабу та єдина система координат}
\label{sec:scale}

Кадр і ортофотоплан мають різну просторову роздільність: один піксель кадру відповідає приблизно $s_f = 0{,}27$ пікселя карти (значення визначено калібруванням, див. п.~\ref{sec:exp-pair}). Якщо подати в метод зіставлення зображення з різним масштабом деталей, частина відповідностей втрачається. Тому перед зіставленням обидва зображення приводяться до спільної роздільності:
\begin{itemize}
    \item фрагмент карти масштабується з коефіцієнтом $k$ (у роботі $k = 1$, тобто фрагмент подається як $512 \times 512$);
    \item кадр масштабується з коефіцієнтом $s_f \cdot k$: кадр $1920 \times 1080$ перетворюється на $\approx 518 \times 292$ пікселі. Кадр використовується повністю, без обрізання.
\end{itemize}

Методи зіставлення повертають координати відповідностей у пікселях зменшених (<<робочих>>) зображень. Система переносить їх назад в оригінальні пікселі кадру та фрагмента з урахуванням узгодження центрів пікселів, яке використовується під час масштабування:
\begin{equation}
x_{\text{ориг}} = \frac{x_{\text{роб}} + 0{,}5}{s_x} - 0{,}5, \qquad
y_{\text{ориг}} = \frac{y_{\text{роб}} + 0{,}5}{s_y} - 0{,}5,
\label{eq:backmap}
\end{equation}
де $s_x, s_y$ --- фактичні коефіцієнти масштабування за осями (вони можуть дещо відрізнятися від заданого через округлення розмірів). Унаслідок цього гомографія $H$ оцінюється між оригінальними пікселями кадру та оригінальними пікселями фрагмента. Усі геометричні величини (поріг RANSAC, похибка репроєкції, координати центра й контуру кадру) виражаються в пікселях карти і не залежать ні від методу зіставлення, ні від робочого масштабу.

\subsection{Реалізація методів зіставлення}

\textbf{Спільний інтерфейс.}\\
Обидва методи реалізовано як класи зі спільним інтерфейсом, що складається з двох кроків:
\begin{enumerate}
    \item \code{extract(image)} --- підготовка одного зображення: для LightGlue --- виділення ключових точок і дескрипторів SuperPoint; для LoFTR, який не має окремого етапу детекції, --- лише формування тензора;
    \item \code{match\_features(f\textsubscript{0}, f\textsubscript{1})} --- встановлення відповідностей між двома підготовленими зображеннями.
\end{enumerate}
Таке розділення дає змогу обчислювати ознаки кадру один раз під час перебору багатьох фрагментів, а ознаки фрагментів карти --- заздалегідь. Решта системи працює лише з цим інтерфейсом і не залежить від конкретного методу.

\textbf{LightGlue + SuperPoint.}\\
SuperPoint виділяє до 2048 ключових точок на зображення; LightGlue використовується з параметрами авторів: поріг адаптивної зупинки 0,95, поріг відсікання точок 0,99, поріг прийняття відповідності 0,1. Під час аналізу прототипу курсової роботи виявлено, що функція виділення ознак SuperPoint за замовчуванням масштабує вхідне зображення до 1024~пікселів за довшою стороною. Унаслідок цього в курсовій роботі LightGlue фактично працював із зображеннями $1024 \times 1024$, а LoFTR --- з $512 \times 512$, тобто методи порівнювалися за різних умов. У розробленій системі це неявне масштабування вимкнено, і масштаб для обох методів задається однаково (п.~\ref{sec:scale}).

\textbf{LoFTR.}\\
Використано реалізацію бібліотеки Kornia з вагами \textit{outdoor}; зображення перетворюються у відтінки сірого. Грубий рівень LoFTR працює на сітці з кроком 8 пікселів, тому розміри вхідних зображень округлюються до кратних 8; відповідні фактичні коефіцієнти масштабу враховуються у формулі~(\ref{eq:backmap}).

\textbf{Кешування ознак фрагментів карти.}\\
Для LightGlue ознаки SuperPoint усіх 384 фрагментів обчислюються один раз і зберігаються на диску (координати точок, оцінки та дескриптори у форматі float32, без втрати точності): $523$--$1609$ точок на фрагмент, загальний обсяг $\approx 347$~МБ, час побудови $\approx 14$~с. Разом із кешем зберігаються параметри, з якими його побудовано (параметри SuperPoint, робочий масштаб, контрольні суми таблиці фрагментів і карти); кеш, що не відповідає поточним параметрам, система відмовляється використовувати. Для LoFTR кешування неможливе, оскільки його ознаки залежать від обох зображень пари.

\textbf{Вимірювання часу.}\\
Перед вимірюваннями виконується <<прогрівання>> моделей на зображеннях тих самих розмірів; для врахування асинхронних обчислень на GPU кожен вимірюваний блок обрамлюється синхронізацією CUDA. Час розділено на етапи: обробка кадру, обробка фрагмента, зіставлення, геометрична перевірка. Читання зображень із диска до вимірюваного часу не входить.

\subsection{Геометрична перевірка відповідностей}

\textbf{Оцінювання гомографії.}\\
За отриманими відповідностями $\{(\mathbf{p}_i, \mathbf{q}_i)\}$ (точка кадру --- точка фрагмента) оцінюється гомографія $H$ (див. розділ~1). Використовується алгоритм RANSAC \cite{ransac} із порогом $\tau = 3$~\px{}, максимальною кількістю ітерацій 5000 і довірчою ймовірністю 0,999; для відтворюваності генератор випадкових чисел ініціалізується фіксованим значенням. Як альтернативний оцінювач реалізовано MAGSAC++ \cite{magsac}, який замість жорсткого порогу зважує відповідності за їхньою похибкою; порівняння оцінювачів наведено в розділі~\ref{sec:experiments}.

Відповідність вважається \emph{геометрично узгодженою}, якщо пряма похибка репроєкції не перевищує порогу:
\begin{equation}
e_i = \left\| \pi(H \tilde{\mathbf{p}}_i) - \mathbf{q}_i \right\|_2 \le \tau,
\end{equation}
де $\tilde{\mathbf{p}}_i$ --- однорідні координати точки, $\pi([u, v, w]^T) = (u/w,\, v/w)$ --- перехід від однорідних координат до декартових. Якщо кількість відповідностей менша за мінімально необхідну для успіху (15), RANSAC не запускається: на наборі, що складається лише з хибних пар, \code{cv2.RANSAC} відпрацьовує всі 5000 ітерацій ($\approx 0{,}1$~с) без шансу на успішний результат.

\textbf{Перевірка правдоподібності гомографії.}\\
RANSAC може знайти формально узгоджену, але фізично неможливу модель (наприклад, на повторюваних структурах). Тому для знайденої $H$ обчислюється якобіан $J$ відображення в центрі кадру $\mathbf{c}$ --- локальне афінне наближення гомографії --- і з нього визначаються:
\begin{align}
\sigma &= \sqrt{|\det J|} && \text{--- масштаб (пікселів карти на піксель кадру)}, \\
\theta &= \operatorname{atan2}(J_{21} - J_{12},\; J_{11} + J_{22}) && \text{--- кут повороту кадру відносно карти}, \\
\kappa &= \sigma_{\max}(J) / \sigma_{\min}(J) && \text{--- анізотропія (відношення сингулярних чисел)}.
\end{align}
Гомографія визнається правдоподібною, якщо виконуються всі умови:
\begin{enumerate}
    \item третя однорідна координата має однаковий знак в усіх кутах кадру (рамка кадру не перетинає <<лінію нескінченності>>);
    \item $\det J > 0$ --- немає дзеркального відображення;
    \item проєкція рамки кадру є опуклим чотирикутником;
    \item $\kappa \le 1{,}5$ --- для надирної камери масштаб за осями майже однаковий;
    \item $s_f / 2 \le \sigma \le 2 s_f$ --- масштаб узгоджується з очікуваним.
\end{enumerate}
Кут $\theta$ надалі використовується як оцінка курсу БпЛА відносно карти, а масштаб $\sigma$ --- як непрямий показник висоти польоту.

\textbf{Критерій успішного зіставлення пари.}\\
Зіставлення пари <<кадр --- фрагмент>> вважається успішним, якщо гомографію знайдено, кількість геометрично узгоджених відповідностей $N_{\text{in}} \ge 15$ і гомографія правдоподібна. Інакше система фіксує конкретну причину невдачі.

\subsection{Автоматичний пошук фрагмента карти}
\label{sec:search}

Для кадру, положення якого невідоме, конвеєр зіставлення виконується з кожним фрагментом пулу (усієї карти, заданої прямокутної області або прорідженої сітки фрагментів). Ознаки кадру обчислюються один раз. Вибір найкращого фрагмента виконується у два кроки.

\textbf{Крок 1: жорсткий фільтр.}
Фрагмент-кандидат відкидається, якщо не виконується хоча б одна умова:
\begin{enumerate}
    \item гомографію знайдено і вона правдоподібна;
    \item $N_{\text{in}} \ge 15$;
    \item частка геометрично узгоджених відповідностей $r = N_{\text{in}} / N_{\text{all}} \ge 0{,}15$;
    \item середня похибка репроєкції $\bar{e} \le 3{,}5$~\px{};
    \item проєкція центра кадру $(x_{\text{лок}}, y_{\text{лок}})$ потрапляє в межі фрагмента з допуском $m = 30$~\px{}: $-m \le x_{\text{лок}} \le S - 1 + m$ (аналогічно для $y$), тобто фрагмент є основним носієм центра кадру;
    \item проєкція центра кадру лежить у межах карти.
\end{enumerate}
Умова~4 за порогу RANSAC $\tau = 3$~\px{} виконується завжди і є запобіжником на випадок зміни порогу.

\textbf{Крок 2: бал ранжування.}
Серед кандидатів, що пройшли фільтр, обирається фрагмент із найбільшим значенням
\begin{equation}
\operatorname{Score} = N_{\text{in}} \cdot \frac{N_{\text{in}}}{N_{\text{all}}} \cdot \exp\!\left(-\frac{\bar{e}}{\sigma_r}\right), \qquad \sigma_r = 2~\text{\px{}}.
\label{eq:score}
\end{equation}
Бал зростає з кількістю та часткою геометрично узгоджених відповідностей і плавно зменшується з похибкою репроєкції. На відміну від відношення $N_{\text{in}} / \bar{e}$, він не стає необмежено великим за малої похибки на кількох випадкових точках.

\textbf{Обчислення координат на карті.}\\
Центр кадру $\mathbf{c} = ((W_f - 1)/2,\, (H_f - 1)/2)$ (у конвенції центрів пікселів) переноситься гомографією обраного фрагмента з явною дегомогенізацією:
\begin{equation}
\tilde{\mathbf{p}} = H \begin{bmatrix} c_x \\ c_y \\ 1 \end{bmatrix}, \qquad
x_{\text{лок}} = \frac{\tilde{p}_1}{\tilde{p}_3}, \quad y_{\text{лок}} = \frac{\tilde{p}_2}{\tilde{p}_3},
\end{equation}
після чого обчислюються глобальні координати на карті:
\begin{equation}
X = X_0 + x_{\text{лок}}, \qquad Y = Y_0 + y_{\text{лок}},
\end{equation}
де $(X_0, Y_0)$ --- координати кута фрагмента з таблиці метаданих. Аналогічно переносяться чотири кути кадру, що дає контур огляду камери на карті. Точки, для яких $|\tilde{p}_3| \approx 0$, вважаються невизначеними й відкидаються.

\subsection{Локалізація послідовності кадрів}
\label{sec:tracking}

Повний перебір 384 фрагментів для кожного кадру є надлишковим: за 1~с апарат зміщується лише на кілька десятків пікселів карти. Тому для послідовності кадрів реалізовано трекер зі станами \textsc{lost} та \textsc{tracking}.

\textbf{Стан \textsc{lost}} (перший кадр або втрата супроводу): виконується глобальний пошук (п.~\ref{sec:search}). Якщо переможець має бал $\operatorname{Score} \ge 5$, положення приймається (статус \textsc{initialized}) і трекер переходить у стан \textsc{tracking}.

\textbf{Стан \textsc{tracking}}: пошук виконується лише серед фрагментів вікна $3 \times 3$ сітки (не більше 9 фрагментів) навколо базового фрагмента. Базовим обирається фрагмент, центр якого найближчий до \emph{прогнозованого} положення. Прогноз будується за двома останніми прийнятими положеннями $(t_{k-2}, P_{k-2})$, $(t_{k-1}, P_{k-1})$ у припущенні сталої швидкості:
\begin{equation}
\hat{P}_k = P_{k-1} + \mathbf{v} \cdot (t_k - t_{k-1}), \qquad
\mathbf{v} = \frac{P_{k-1} - P_{k-2}}{t_{k-1} - t_{k-2}},
\label{eq:extrapolation}
\end{equation}
де модуль швидкості обмежується значенням $V_{\max}$. За сталого інтервалу між кадрами це збігається з $\hat{P}_k = P_{k-1} + (P_{k-1} - P_{k-2})$.

\textbf{Кінематична перевірка.}\\
Нехай $P_{\text{ост}}$ --- останнє \emph{прийняте} положення, отримане в момент $t_{\text{ост}}$. Нова оцінка $P_k$ приймається, якщо
\begin{equation}
\| P_k - P_{\text{ост}} \|_2 \le D_{\max} = V_{\max} \cdot (t_k - t_{\text{ост}}) + D_{\text{запас}},
\end{equation}
де $V_{\max} = 120$~\px{}/с, $D_{\text{запас}} = 30$~\px{}; за частоти 1~кадр/с $D_{\max} = 150$~\px{}. Порівняння саме з останнім прийнятим положенням означає, що після кількох пропущених кадрів допустимий радіус зростає пропорційно часу. Серед кандидатів, що пройшли фільтр п.~\ref{sec:search}, обирається найкращий за балом, який проходить і кінематичну перевірку. Якщо такого кандидата немає, положення кадру екстраполюється за формулою~(\ref{eq:extrapolation}) (статус \textsc{interpolated}), а відкинута оцінка зберігається для аналізу. Після трьох поспіль кадрів без прийнятого положення трекер переходить у стан \textsc{lost}.

\textbf{Компенсація курсу.}\\
SuperPoint/LightGlue і LoFTR не інваріантні до великих поворотів зображення (кількісне підтвердження наведено в п.~\ref{sec:exp-robust}), а під час розвороту БпЛА кадр повертається відносно карти. Кут <<кадр $\to$ карта>> відомий із гомографій прийнятих кадрів ($\theta$, п.~\ref{sec:search}), тому перед зіставленням кадр повертається на прогнозований кут $\hat{\theta}_k$, і метод зіставлення отримує кадр, орієнтований так само, як карта. Прогноз кута будується аналогічно до~(\ref{eq:extrapolation}) з урахуванням переходу через $\pm 180^\circ$, а кутова швидкість обмежується значенням $45^\circ$/с. Полотно під час повороту розширюється, а відповідності, що потрапили на фон поза кадром, відкидаються. Перетворення <<оригінальний кадр $\to$ робоче зображення>> є композицією масштабування та повороту $T = R(\hat{\theta}_k) \cdot S$, і координати відповідностей переносяться назад точним оберненим перетворенням $T^{-1}$. Тому гомографія, як і раніше, оцінюється в оригінальних пікселях кадру.

Якщо з прогнозованим курсом фрагмент не знайдено, у стані \textsc{tracking} послідовно перевіряються гіпотези $\hat{\theta}_k \pm 45^\circ$ і $\hat{\theta}_k \pm 90^\circ$. У стані \textsc{lost}, де повний перебір карти дорогий, кожен наступний глобальний пошук використовує чергову гіпотезу з циклу $0^\circ, 90^\circ, 180^\circ, -90^\circ$ відносно останнього відомого курсу.

Узагальнено обробка одного кадру в стані \textsc{tracking} складається з таких кроків:
\begin{enumerate}
    \item обчислити прогноз положення $\hat{P}_k$ і курсу $\hat{\theta}_k$;
    \item сформувати пул із не більше ніж 9 фрагментів навколо фрагмента, найближчого до $\hat{P}_k$;
    \item для кожної гіпотези курсу, починаючи з $\hat{\theta}_k$:
    \begin{enumerate}
        \item повернути та масштабувати кадр, обчислити його ознаки;
        \item зіставити кадр із фрагментами пулу, виконати геометричну перевірку, фільтр і ранжування;
        \item якщо найкращий кандидат проходить кінематичну перевірку --- прийняти положення (статус \textsc{tracking}) і завершити обробку кадру;
    \end{enumerate}
    \item якщо жодна гіпотеза не дала прийнятного положення --- екстраполювати положення (статус \textsc{interpolated}); після трьох таких кадрів поспіль перейти в стан \textsc{lost}.
\end{enumerate}

\subsection{Модульна структура та відтворюваність}

Програмну систему організовано як пакет із незалежними модулями (табл.~\ref{tab:modules}) і набір сценаріїв командного рядка для кожного етапу.

\begin{table}[H]
\centering
\caption{Основні модулі програмної системи}
\label{tab:modules}
\small
\begin{tabular}{p{0.36\linewidth} p{0.56\linewidth}}
\toprule
\textbf{Модуль} & \textbf{Призначення} \\
\midrule
\code{preprocessing/tiles.py}, \code{frames.py} & розбиття карти, вилучення кадрів, метадані \\
\code{preprocessing/images.py} & читання/запис, масштабування, поворот, перенесення координат, спотворення для експериментів \\
\code{preprocessing/cache.py} & дисковий кеш ознак фрагментів \\
\code{matching/lightglue.py}, \code{loftr.py} & реалізації методів зі спільним інтерфейсом \\
\code{matching/pipeline.py} & єдиний конвеєр <<кадр + фрагмент $\to$ відповідності $\to$ RANSAC $\to$ метрики>> \\
\code{geometry/homography.py} & RANSAC/MAGSAC++, похибка репроєкції, перевірка правдоподібності \\
\code{evaluation/scoring.py} & фільтр і бал ранжування фрагментів \\
\code{localization/coordinates.py} & перенесення точок на карту, контур огляду \\
\code{localization/search.py}, \code{tracker.py} & пошук фрагмента, трекер послідовності \\
\code{evaluation/metrics.py}, \code{ground\_truth.py} & метрики, еталонні координати \\
\code{visualization/} & рисунки відповідностей, положення та траєкторій \\
\bottomrule
\end{tabular}
\end{table}

Для відтворюваності результатів кожен запуск записує, разом із результатами, файл метаданих з усіма параметрами конвеєра й методу зіставлення, ідентифікатором версії коду та відомостями про обчислювальне середовище. Генератори випадкових чисел RANSAC і спотворень ініціалізуються фіксованими значеннями. Коректність перенесення координат, оцінювання гомографії, критеріїв фільтрації та логіки трекера перевіряється набором із 72 автоматичних тестів на синтетичних даних, для яких точний результат відомий заздалегідь; тести не потребують нейромереж і виконуються за кілька секунд.

\subsection{Візуалізація результатів}

Система формує такі графічні матеріали:
\begin{itemize}
    \item відповідності між кадром і фрагментом карти: геометрично узгоджені --- зеленим, решта --- червоним;
    \item контур огляду камери і центр кадру на ортофотоплані разом із сіткою перевірених фрагментів і обраним фрагментом;
    \item траєкторію польоту на карті з позначенням статусів кадрів і напрямку руху;
    \item графіки похибок, часу обробки та результатів експериментів (розділ~\ref{sec:experiments}).
\end{itemize}

% Потрібно в преамбулі: \usetikzlibrary{arrows.meta,positioning,babel}
% =============================================================================
\providecommand{\px}{пкс}               % пікселі карти
\providecommand{\code}[1]{\texttt{#1}}  % імена файлів, модулів і функцій
\newpage
\section{Програмне забезпечення та програмна реалізація системи}
\label{sec:implementation}

\subsection{Постановка задачі та загальна архітектура системи}

Розроблена система розв'язує задачу абсолютної візуальної локалізації БпЛА: за кадром із бортової камери та ортофотопланом місцевості визначити, над якою точкою карти перебуває апарат, а за послідовністю кадрів --- відновити траєкторію польоту. Інформація про те, якому фрагменту карти відповідає кадр, системі не надається: відповідну ділянку ортофотоплану система знаходить автоматично.

Вхідними даними є:
\begin{itemize}
    \item ортофотоплан місцевості розміром $W \times H$ пікселів;
    \item відеозапис польоту БпЛА або послідовність кадрів із часовими мітками.
\end{itemize}

Вихідними даними є:
\begin{itemize}
    \item для кожного кадру --- координати $(X, Y)$ проєкції центра кадру на ортофотоплан і контур огляду камери (чотирикутник на карті);
    \item для послідовності кадрів --- траєкторія польоту з позначкою статусу кожного кадру та показниками якості зіставлення.
\end{itemize}

Робота системи ґрунтується на таких припущеннях:
\begin{itemize}
    \item камера спрямована приблизно вертикально вниз (надирна зйомка);
    \item сцену в межах кадру можна наближено вважати плоскою, тому перетворення <<кадр $\to$ карта>> описується гомографією \cite{hartley};
    \item відношення просторової роздільності кадру та карти відоме з точністю до множника $\approx 2$ (див. п.~\ref{sec:scale});
    \item ортофотоплан не має географічної прив'язки, тому всі координати та похибки обчислюються в пікселях карти (\px{}), без переведення в метри.
\end{itemize}

Загальну схему обробки наведено на рисунку~\ref{fig:pipeline}. Система складається з офлайн-етапу підготовки карти, який виконується один раз, та онлайн-етапу обробки кадрів. Ключовим елементом онлайн-етапу є єдиний конвеєр зіставлення пари <<кадр --- фрагмент карти>>, який однаково застосовується в усіх режимах роботи: під час перевірки відомої пари, автоматичного пошуку фрагмента та трекінгу послідовності.

\begin{figure}[H]
\centering
\begin{tikzpicture}[
    node distance=5mm and 6mm,
    box/.style={draw, rounded corners=2pt, align=center, font=\small,
                minimum width=42mm, minimum height=9mm, inner sep=3pt},
    off/.style={box, fill=gray!12},
    on/.style={box, fill=blue!7},
    dec/.style={box, fill=orange!12},
    arr/.style={-{Stealth[length=2.2mm]}, thick}]
  % офлайн-етап
  \node[off] (map) {Ортофотоплан\\$6200\times4320$};
  \node[off, below=of map] (tiles) {Розбиття на фрагменти\\$512\times512$, перекриття 50\,\%};
  \node[off, below=of tiles] (cache) {Кеш ознак фрагментів\\(SuperPoint)};
  % онлайн-етап
  \node[on, right=16mm of map] (frame) {Кадр БпЛА\\$1920\times1080$};
  \node[on, below=of frame] (prep) {Нормалізація масштабу\\і компенсація курсу};
  \node[dec, below=of prep] (pool) {Пул фрагментів:\\вся карта / вікно $3\times3$};
  \node[on, below=of pool] (match) {Зіставлення\\LightGlue / LoFTR};
  \node[on, below=of match] (geom) {RANSAC + гомографія,\\перевірка правдоподібності};
  \node[dec, below=of geom] (rank) {Фільтр і ранжування\\фрагментів};
  \node[on, below=of rank] (coord) {Координати $(X,Y)$\\і контур огляду};
  \node[dec, below=of coord] (track) {Кінематична перевірка,\\траєкторія};
  \draw[arr] (map) -- (tiles);
  \draw[arr] (tiles) -- (cache);
  \draw[arr] (frame) -- (prep);
  \draw[arr] (prep) -- (pool);
  \draw[arr] (pool) -- (match);
  \draw[arr] (match) -- (geom);
  \draw[arr] (geom) -- (rank);
  \draw[arr] (rank) -- (coord);
  \draw[arr] (coord) -- (track);
  \draw[arr] (cache.south) |- (match.west);
  \draw[arr, dashed] (track.east) -- ++(6mm,0) |- node[pos=0.25, right, font=\footnotesize, align=left]{прогноз\\положення\\і курсу} (pool.east);
\end{tikzpicture}
\caption{Загальна схема роботи системи візуальної локалізації (сірим позначено офлайн-етап)}
\label{fig:pipeline}
\end{figure}

\subsection{Використані програмні засоби}

Систему реалізовано мовою Python~3.11. Використано такі бібліотеки:
\begin{itemize}[label=\textbf{•}]
    \item \textbf{PyTorch} 2.14 з підтримкою CUDA~12.6 \cite{pytorch} --- виконання нейромережевих моделей на GPU;
    \item \textbf{LightGlue} (офіційна реалізація авторів) --- моделі SuperPoint \cite{superpoint} і LightGlue \cite{lightglue};
    \item \textbf{Kornia} 0.8 \cite{kornia} --- реалізація LoFTR \cite{loftr} з попередньо навченими вагами \textit{outdoor};
    \item \textbf{OpenCV} 5.0 \cite{opencv} --- читання та масштабування зображень, оцінювання гомографії методами RANSAC \cite{ransac} і MAGSAC++ \cite{magsac}, геометричні перетворення;
    \item \textbf{NumPy} і \textbf{pandas} --- числові обчислення, робота з таблицями метаданих і результатів;
    \item \textbf{Matplotlib} --- побудова графіків;
    \item \textbf{pytest} --- автоматичне тестування модулів.
\end{itemize}

Нейромережі в роботі не навчалися: використано опубліковані попередньо навчені ваги моделей.

\subsection{Тестове середовище}

Розроблення та експерименти виконувалися на ноутбуці з такими характеристиками:
\begin{itemize}[label=\textbf{--}]
    \item \textbf{процесор:} AMD Ryzen 5 4600H;
    \item \textbf{оперативна пам'ять:} 8 ГБ;
    \item \textbf{відеокарта:} NVIDIA GeForce GTX 1650 (4 ГБ відеопам'яті);
    \item \textbf{CUDA:} 12.6 (збірка PyTorch \texttt{cu126});
    \item \textbf{операційна система:} Windows 11.
\end{itemize}

\subsection{Вхідні дані}

Для моделювання використано ортофотоплан району розміром $6200 \times 4320$ пікселів (формат RGBA з повністю непрозорим альфа-каналом) і відеозапис польоту БпЛА DJI (файл \code{DJI\_0331.MP4}): роздільність $1920 \times 1080$, частота 29,97~кадр/с, тривалість $\approx 94$~с, кодек HEVC, камера спрямована вниз.

Відео та ортофотоплан знято в різні пори року: ортофотоплан --- влітку, відео --- взимку. Через це рослинність, тіні та покриття ґрунту на кадрах і на карті суттєво відрізняються, що ускладнює зіставлення.

\subsection{Підготовка ортофотоплану}

\textbf{Розбиття на фрагменти з перекриттям.}\\
Нейромережеві методи зіставлення розраховані на зображення розміром порядку сотень пікселів, тому ортофотоплан розбивається на квадратні фрагменти (тайли) розміром $S = 512$ пікселів із часткою перекриття $o = 0{,}5$. Крок сітки визначається як
\begin{equation}
s = \lfloor S \cdot (1 - o) \rfloor = 256 .
\end{equation}
Координати лівих верхніх кутів фрагментів уздовж осі $x$ утворюють послідовність $0, s, 2s, \dots$; якщо регулярна сітка не доходить до правого краю карти, додається ще один крайовий фрагмент із початком у $W - S$ (аналогічно для осі $y$). Завдяки цьому всі фрагменти мають однаковий розмір і кожен піксель карти покрито. Кількість фрагментів:
\begin{equation}
N = \left( \left\lceil \frac{W - S}{s} \right\rceil + 1 \right) \cdot \left( \left\lceil \frac{H - S}{s} \right\rceil + 1 \right) = 24 \cdot 16 = 384 .
\end{equation}
Для кожного фрагмента в таблицю метаданих \code{tiles.csv} записуються ідентифікатор і координати $(X_0, Y_0)$ лівого верхнього кута в пікселях повної карти. Ці координати є основою для переходу від локальних координат фрагмента до глобальних координат карти.

Перекриття 50\,\% гарантує, що будь-яка ділянка карти, менша за половину фрагмента, повністю потрапляє принаймні в один фрагмент. Водночас, як показано в розділі~\ref{sec:experiments}, контур огляду камери на карті має розмір $\approx 518 \times 292$~\px{} і є ширшим за фрагмент, тому кадр, як правило, лише частково перекривається з кожним окремим фрагментом.

\textbf{Перевірка цілісності даних.}\\
Після розбиття автоматично перевіряється, що кожен фрагмент попіксельно збігається з відповідною ділянкою вихідної карти, усі фрагменти лежать у межах карти, їх кількість відповідає формулі, а таблиця метаданих не містить пропусків і дублікатів.

\subsection{Вилучення кадрів із відео}

З відео вибирається один кадр на секунду: кожен $n$-й кадр, де $n = \operatorname{round}(f_{\text{відео}} / f_{\text{ціль}}) = \operatorname{round}(29{,}97 / 1) = 30$. Часова мітка кадру обчислюється як $t = i / f_{\text{відео}}$, де $i$ --- номер кадру у відео. Кадри декодуються послідовно (без позиціювання за номером, яке ненадійне для кодеків H.264/HEVC) і зберігаються без втрат у форматі PNG. У результаті отримано 95 кадрів із кроком $\approx 1{,}001$~с.

\subsection{Нормалізація масштабу та єдина система координат}
\label{sec:scale}

Кадр і ортофотоплан мають різну просторову роздільність: один піксель кадру відповідає приблизно $s_f = 0{,}27$ пікселя карти (значення визначено калібруванням, див. п.~\ref{sec:exp-pair}). Якщо подати в метод зіставлення зображення з різним масштабом деталей, частина відповідностей втрачається. Тому перед зіставленням обидва зображення приводяться до спільної роздільності:
\begin{itemize}
    \item фрагмент карти масштабується з коефіцієнтом $k$ (у роботі $k = 1$, тобто фрагмент подається як $512 \times 512$);
    \item кадр масштабується з коефіцієнтом $s_f \cdot k$: кадр $1920 \times 1080$ перетворюється на $\approx 518 \times 292$ пікселі. Кадр використовується повністю, без обрізання.
\end{itemize}

Методи зіставлення повертають координати відповідностей у пікселях зменшених (<<робочих>>) зображень. Система переносить їх назад в оригінальні пікселі кадру та фрагмента з урахуванням узгодження центрів пікселів, яке використовується під час масштабування:
\begin{equation}
x_{\text{ориг}} = \frac{x_{\text{роб}} + 0{,}5}{s_x} - 0{,}5, \qquad
y_{\text{ориг}} = \frac{y_{\text{роб}} + 0{,}5}{s_y} - 0{,}5,
\label{eq:backmap}
\end{equation}
де $s_x, s_y$ --- фактичні коефіцієнти масштабування за осями (вони можуть дещо відрізнятися від заданого через округлення розмірів). Унаслідок цього гомографія $H$ оцінюється між оригінальними пікселями кадру та оригінальними пікселями фрагмента. Усі геометричні величини (поріг RANSAC, похибка репроєкції, координати центра й контуру кадру) виражаються в пікселях карти і не залежать ні від методу зіставлення, ні від робочого масштабу.

\subsection{Реалізація методів зіставлення}

\textbf{Спільний інтерфейс.}\\
Обидва методи реалізовано як класи зі спільним інтерфейсом, що складається з двох кроків:
\begin{enumerate}
    \item \code{extract(image)} --- підготовка одного зображення: для LightGlue --- виділення ключових точок і дескрипторів SuperPoint; для LoFTR, який не має окремого етапу детекції, --- лише формування тензора;
    \item \code{match\_features(f\textsubscript{0}, f\textsubscript{1})} --- встановлення відповідностей між двома підготовленими зображеннями.
\end{enumerate}
Таке розділення дає змогу обчислювати ознаки кадру один раз під час перебору багатьох фрагментів, а ознаки фрагментів карти --- заздалегідь. Решта системи працює лише з цим інтерфейсом і не залежить від конкретного методу.

\textbf{LightGlue + SuperPoint.}\\
SuperPoint виділяє до 2048 ключових точок на зображення; LightGlue використовується з параметрами авторів: поріг адаптивної зупинки 0,95, поріг відсікання точок 0,99, поріг прийняття відповідності 0,1. Під час аналізу прототипу курсової роботи виявлено, що функція виділення ознак SuperPoint за замовчуванням масштабує вхідне зображення до 1024~пікселів за довшою стороною. Унаслідок цього в курсовій роботі LightGlue фактично працював із зображеннями $1024 \times 1024$, а LoFTR --- з $512 \times 512$, тобто методи порівнювалися за різних умов. У розробленій системі це неявне масштабування вимкнено, і масштаб для обох методів задається однаково (п.~\ref{sec:scale}).

\textbf{LoFTR.}\\
Використано реалізацію бібліотеки Kornia з вагами \textit{outdoor}; зображення перетворюються у відтінки сірого. Грубий рівень LoFTR працює на сітці з кроком 8 пікселів, тому розміри вхідних зображень округлюються до кратних 8; відповідні фактичні коефіцієнти масштабу враховуються у формулі~(\ref{eq:backmap}).

\textbf{Кешування ознак фрагментів карти.}\\
Для LightGlue ознаки SuperPoint усіх 384 фрагментів обчислюються один раз і зберігаються на диску (координати точок, оцінки та дескриптори у форматі float32, без втрати точності): $523$--$1609$ точок на фрагмент, загальний обсяг $\approx 347$~МБ, час побудови $\approx 14$~с. Разом із кешем зберігаються параметри, з якими його побудовано (параметри SuperPoint, робочий масштаб, контрольні суми таблиці фрагментів і карти); кеш, що не відповідає поточним параметрам, система відмовляється використовувати. Для LoFTR кешування неможливе, оскільки його ознаки залежать від обох зображень пари.

\textbf{Вимірювання часу.}\\
Перед вимірюваннями виконується <<прогрівання>> моделей на зображеннях тих самих розмірів; для врахування асинхронних обчислень на GPU кожен вимірюваний блок обрамлюється синхронізацією CUDA. Час розділено на етапи: обробка кадру, обробка фрагмента, зіставлення, геометрична перевірка. Читання зображень із диска до вимірюваного часу не входить.

\subsection{Геометрична перевірка відповідностей}

\textbf{Оцінювання гомографії.}\\
За отриманими відповідностями $\{(\mathbf{p}_i, \mathbf{q}_i)\}$ (точка кадру --- точка фрагмента) оцінюється гомографія $H$ (див. розділ~1). Використовується алгоритм RANSAC \cite{ransac} із порогом $\tau = 3$~\px{}, максимальною кількістю ітерацій 5000 і довірчою ймовірністю 0,999; для відтворюваності генератор випадкових чисел ініціалізується фіксованим значенням. Як альтернативний оцінювач реалізовано MAGSAC++ \cite{magsac}, який замість жорсткого порогу зважує відповідності за їхньою похибкою; порівняння оцінювачів наведено в розділі~\ref{sec:experiments}.

Відповідність вважається \emph{геометрично узгодженою}, якщо пряма похибка репроєкції не перевищує порогу:
\begin{equation}
e_i = \left\| \pi(H \tilde{\mathbf{p}}_i) - \mathbf{q}_i \right\|_2 \le \tau,
\end{equation}
де $\tilde{\mathbf{p}}_i$ --- однорідні координати точки, $\pi([u, v, w]^T) = (u/w,\, v/w)$ --- перехід від однорідних координат до декартових. Якщо кількість відповідностей менша за мінімально необхідну для успіху (15), RANSAC не запускається: на наборі, що складається лише з хибних пар, \code{cv2.RANSAC} відпрацьовує всі 5000 ітерацій ($\approx 0{,}1$~с) без шансу на успішний результат.

\textbf{Перевірка правдоподібності гомографії.}\\
RANSAC може знайти формально узгоджену, але фізично неможливу модель (наприклад, на повторюваних структурах). Тому для знайденої $H$ обчислюється якобіан $J$ відображення в центрі кадру $\mathbf{c}$ --- локальне афінне наближення гомографії --- і з нього визначаються:
\begin{align}
\sigma &= \sqrt{|\det J|} && \text{--- масштаб (пікселів карти на піксель кадру)}, \\
\theta &= \operatorname{atan2}(J_{21} - J_{12},\; J_{11} + J_{22}) && \text{--- кут повороту кадру відносно карти}, \\
\kappa &= \sigma_{\max}(J) / \sigma_{\min}(J) && \text{--- анізотропія (відношення сингулярних чисел)}.
\end{align}
Гомографія визнається правдоподібною, якщо виконуються всі умови:
\begin{enumerate}
    \item третя однорідна координата має однаковий знак в усіх кутах кадру (рамка кадру не перетинає <<лінію нескінченності>>);
    \item $\det J > 0$ --- немає дзеркального відображення;
    \item проєкція рамки кадру є опуклим чотирикутником;
    \item $\kappa \le 1{,}5$ --- для надирної камери масштаб за осями майже однаковий;
    \item $s_f / 2 \le \sigma \le 2 s_f$ --- масштаб узгоджується з очікуваним.
\end{enumerate}
Кут $\theta$ надалі використовується як оцінка курсу БпЛА відносно карти, а масштаб $\sigma$ --- як непрямий показник висоти польоту.

\textbf{Критерій успішного зіставлення пари.}\\
Зіставлення пари <<кадр --- фрагмент>> вважається успішним, якщо гомографію знайдено, кількість геометрично узгоджених відповідностей $N_{\text{in}} \ge 15$ і гомографія правдоподібна. Інакше система фіксує конкретну причину невдачі.

\subsection{Автоматичний пошук фрагмента карти}
\label{sec:search}

Для кадру, положення якого невідоме, конвеєр зіставлення виконується з кожним фрагментом пулу (усієї карти, заданої прямокутної області або прорідженої сітки фрагментів). Ознаки кадру обчислюються один раз. Вибір найкращого фрагмента виконується у два кроки.

\textbf{Крок 1: жорсткий фільтр.}
Фрагмент-кандидат відкидається, якщо не виконується хоча б одна умова:
\begin{enumerate}
    \item гомографію знайдено і вона правдоподібна;
    \item $N_{\text{in}} \ge 15$;
    \item частка геометрично узгоджених відповідностей $r = N_{\text{in}} / N_{\text{all}} \ge 0{,}15$;
    \item середня похибка репроєкції $\bar{e} \le 3{,}5$~\px{};
    \item проєкція центра кадру $(x_{\text{лок}}, y_{\text{лок}})$ потрапляє в межі фрагмента з допуском $m = 30$~\px{}: $-m \le x_{\text{лок}} \le S - 1 + m$ (аналогічно для $y$), тобто фрагмент є основним носієм центра кадру;
    \item проєкція центра кадру лежить у межах карти.
\end{enumerate}
Умова~4 за порогу RANSAC $\tau = 3$~\px{} виконується завжди і є запобіжником на випадок зміни порогу.

\textbf{Крок 2: бал ранжування.}
Серед кандидатів, що пройшли фільтр, обирається фрагмент із найбільшим значенням
\begin{equation}
\operatorname{Score} = N_{\text{in}} \cdot \frac{N_{\text{in}}}{N_{\text{all}}} \cdot \exp\!\left(-\frac{\bar{e}}{\sigma_r}\right), \qquad \sigma_r = 2~\text{\px{}}.
\label{eq:score}
\end{equation}
Бал зростає з кількістю та часткою геометрично узгоджених відповідностей і плавно зменшується з похибкою репроєкції. На відміну від відношення $N_{\text{in}} / \bar{e}$, він не стає необмежено великим за малої похибки на кількох випадкових точках.

\textbf{Обчислення координат на карті.}\\
Центр кадру $\mathbf{c} = ((W_f - 1)/2,\, (H_f - 1)/2)$ (у конвенції центрів пікселів) переноситься гомографією обраного фрагмента з явною дегомогенізацією:
\begin{equation}
\tilde{\mathbf{p}} = H \begin{bmatrix} c_x \\ c_y \\ 1 \end{bmatrix}, \qquad
x_{\text{лок}} = \frac{\tilde{p}_1}{\tilde{p}_3}, \quad y_{\text{лок}} = \frac{\tilde{p}_2}{\tilde{p}_3},
\end{equation}
після чого обчислюються глобальні координати на карті:
\begin{equation}
X = X_0 + x_{\text{лок}}, \qquad Y = Y_0 + y_{\text{лок}},
\end{equation}
де $(X_0, Y_0)$ --- координати кута фрагмента з таблиці метаданих. Аналогічно переносяться чотири кути кадру, що дає контур огляду камери на карті. Точки, для яких $|\tilde{p}_3| \approx 0$, вважаються невизначеними й відкидаються.

\subsection{Локалізація послідовності кадрів}
\label{sec:tracking}

Повний перебір 384 фрагментів для кожного кадру є надлишковим: за 1~с апарат зміщується лише на кілька десятків пікселів карти. Тому для послідовності кадрів реалізовано трекер зі станами \textsc{lost} та \textsc{tracking}.

\textbf{Стан \textsc{lost}} (перший кадр або втрата супроводу): виконується глобальний пошук (п.~\ref{sec:search}). Якщо переможець має бал $\operatorname{Score} \ge 5$, положення приймається (статус \textsc{initialized}) і трекер переходить у стан \textsc{tracking}.

\textbf{Стан \textsc{tracking}}: пошук виконується лише серед фрагментів вікна $3 \times 3$ сітки (не більше 9 фрагментів) навколо базового фрагмента. Базовим обирається фрагмент, центр якого найближчий до \emph{прогнозованого} положення. Прогноз будується за двома останніми прийнятими положеннями $(t_{k-2}, P_{k-2})$, $(t_{k-1}, P_{k-1})$ у припущенні сталої швидкості:
\begin{equation}
\hat{P}_k = P_{k-1} + \mathbf{v} \cdot (t_k - t_{k-1}), \qquad
\mathbf{v} = \frac{P_{k-1} - P_{k-2}}{t_{k-1} - t_{k-2}},
\label{eq:extrapolation}
\end{equation}
де модуль швидкості обмежується значенням $V_{\max}$. За сталого інтервалу між кадрами це збігається з $\hat{P}_k = P_{k-1} + (P_{k-1} - P_{k-2})$.

\textbf{Кінематична перевірка.}\\
Нехай $P_{\text{ост}}$ --- останнє \emph{прийняте} положення, отримане в момент $t_{\text{ост}}$. Нова оцінка $P_k$ приймається, якщо
\begin{equation}
\| P_k - P_{\text{ост}} \|_2 \le D_{\max} = V_{\max} \cdot (t_k - t_{\text{ост}}) + D_{\text{запас}},
\end{equation}
де $V_{\max} = 120$~\px{}/с, $D_{\text{запас}} = 30$~\px{}; за частоти 1~кадр/с $D_{\max} = 150$~\px{}. Порівняння саме з останнім прийнятим положенням означає, що після кількох пропущених кадрів допустимий радіус зростає пропорційно часу. Серед кандидатів, що пройшли фільтр п.~\ref{sec:search}, обирається найкращий за балом, який проходить і кінематичну перевірку. Якщо такого кандидата немає, положення кадру екстраполюється за формулою~(\ref{eq:extrapolation}) (статус \textsc{interpolated}), а відкинута оцінка зберігається для аналізу. Після трьох поспіль кадрів без прийнятого положення трекер переходить у стан \textsc{lost}.

\textbf{Компенсація курсу.}\\
SuperPoint/LightGlue і LoFTR не інваріантні до великих поворотів зображення (кількісне підтвердження наведено в п.~\ref{sec:exp-robust}), а під час розвороту БпЛА кадр повертається відносно карти. Кут <<кадр $\to$ карта>> відомий із гомографій прийнятих кадрів ($\theta$, п.~\ref{sec:search}), тому перед зіставленням кадр повертається на прогнозований кут $\hat{\theta}_k$, і метод зіставлення отримує кадр, орієнтований так само, як карта. Прогноз кута будується аналогічно до~(\ref{eq:extrapolation}) з урахуванням переходу через $\pm 180^\circ$, а кутова швидкість обмежується значенням $45^\circ$/с. Полотно під час повороту розширюється, а відповідності, що потрапили на фон поза кадром, відкидаються. Перетворення <<оригінальний кадр $\to$ робоче зображення>> є композицією масштабування та повороту $T = R(\hat{\theta}_k) \cdot S$, і координати відповідностей переносяться назад точним оберненим перетворенням $T^{-1}$. Тому гомографія, як і раніше, оцінюється в оригінальних пікселях кадру.

Якщо з прогнозованим курсом фрагмент не знайдено, у стані \textsc{tracking} послідовно перевіряються гіпотези $\hat{\theta}_k \pm 45^\circ$ і $\hat{\theta}_k \pm 90^\circ$. У стані \textsc{lost}, де повний перебір карти дорогий, кожен наступний глобальний пошук використовує чергову гіпотезу з циклу $0^\circ, 90^\circ, 180^\circ, -90^\circ$ відносно останнього відомого курсу.

Узагальнено обробка одного кадру в стані \textsc{tracking} складається з таких кроків:
\begin{enumerate}
    \item обчислити прогноз положення $\hat{P}_k$ і курсу $\hat{\theta}_k$;
    \item сформувати пул із не більше ніж 9 фрагментів навколо фрагмента, найближчого до $\hat{P}_k$;
    \item для кожної гіпотези курсу, починаючи з $\hat{\theta}_k$:
    \begin{enumerate}
        \item повернути та масштабувати кадр, обчислити його ознаки;
        \item зіставити кадр із фрагментами пулу, виконати геометричну перевірку, фільтр і ранжування;
        \item якщо найкращий кандидат проходить кінематичну перевірку --- прийняти положення (статус \textsc{tracking}) і завершити обробку кадру;
    \end{enumerate}
    \item якщо жодна гіпотеза не дала прийнятного положення --- екстраполювати положення (статус \textsc{interpolated}); після трьох таких кадрів поспіль перейти в стан \textsc{lost}.
\end{enumerate}

\subsection{Модульна структура та відтворюваність}

Програмну систему організовано як пакет із незалежними модулями (табл.~\ref{tab:modules}) і набір сценаріїв командного рядка для кожного етапу.

\begin{table}[H]
\centering
\caption{Основні модулі програмної системи}
\label{tab:modules}
\small
\begin{tabular}{p{0.36\linewidth} p{0.56\linewidth}}
\toprule
\textbf{Модуль} & \textbf{Призначення} \\
\midrule
\code{preprocessing/tiles.py}, \code{frames.py} & розбиття карти, вилучення кадрів, метадані \\
\code{preprocessing/images.py} & читання/запис, масштабування, поворот, перенесення координат, спотворення для експериментів \\
\code{preprocessing/cache.py} & дисковий кеш ознак фрагментів \\
\code{matching/lightglue.py}, \code{loftr.py} & реалізації методів зі спільним інтерфейсом \\
\code{matching/pipeline.py} & єдиний конвеєр <<кадр + фрагмент $\to$ відповідності $\to$ RANSAC $\to$ метрики>> \\
\code{geometry/homography.py} & RANSAC/MAGSAC++, похибка репроєкції, перевірка правдоподібності \\
\code{evaluation/scoring.py} & фільтр і бал ранжування фрагментів \\
\code{localization/coordinates.py} & перенесення точок на карту, контур огляду \\
\code{localization/search.py}, \code{tracker.py} & пошук фрагмента, трекер послідовності \\
\code{evaluation/metrics.py}, \code{ground\_truth.py} & метрики, еталонні координати \\
\code{visualization/} & рисунки відповідностей, положення та траєкторій \\
\bottomrule
\end{tabular}
\end{table}

Для відтворюваності результатів кожен запуск записує, разом із результатами, файл метаданих з усіма параметрами конвеєра й методу зіставлення, ідентифікатором версії коду та відомостями про обчислювальне середовище. Генератори випадкових чисел RANSAC і спотворень ініціалізуються фіксованими значеннями. Коректність перенесення координат, оцінювання гомографії, критеріїв фільтрації та логіки трекера перевіряється набором із 72 автоматичних тестів на синтетичних даних, для яких точний результат відомий заздалегідь; тести не потребують нейромереж і виконуються за кілька секунд.

\subsection{Візуалізація результатів}

Система формує такі графічні матеріали:
\begin{itemize}
    \item відповідності між кадром і фрагментом карти: геометрично узгоджені --- зеленим, решта --- червоним;
    \item контур огляду камери і центр кадру на ортофотоплані разом із сіткою перевірених фрагментів і обраним фрагментом;
    \item траєкторію польоту на карті з позначенням статусів кадрів і напрямку руху;
    \item графіки похибок, часу обробки та результатів експериментів (розділ~\ref{sec:experiments}).
\end{itemize}

% Потрібно в преамбулі: \usetikzlibrary{arrows.meta,positioning,babel}
% =============================================================================
\providecommand{\px}{пкс}               % пікселі карти
\providecommand{\code}[1]{\texttt{#1}}  % імена файлів, модулів і функцій
\newpage
\section{Програмне забезпечення та програмна реалізація системи}
\label{sec:implementation}

\subsection{Постановка задачі та загальна архітектура системи}

Розроблена система розв'язує задачу абсолютної візуальної локалізації БпЛА: за кадром із бортової камери та ортофотопланом місцевості визначити, над якою точкою карти перебуває апарат, а за послідовністю кадрів --- відновити траєкторію польоту. Інформація про те, якому фрагменту карти відповідає кадр, системі не надається: відповідну ділянку ортофотоплану система знаходить автоматично.

Вхідними даними є:
\begin{itemize}
    \item ортофотоплан місцевості розміром $W \times H$ пікселів;
    \item відеозапис польоту БпЛА або послідовність кадрів із часовими мітками.
\end{itemize}

Вихідними даними є:
\begin{itemize}
    \item для кожного кадру --- координати $(X, Y)$ проєкції центра кадру на ортофотоплан і контур огляду камери (чотирикутник на карті);
    \item для послідовності кадрів --- траєкторія польоту з позначкою статусу кожного кадру та показниками якості зіставлення.
\end{itemize}

Робота системи ґрунтується на таких припущеннях:
\begin{itemize}
    \item камера спрямована приблизно вертикально вниз (надирна зйомка);
    \item сцену в межах кадру можна наближено вважати плоскою, тому перетворення <<кадр $\to$ карта>> описується гомографією \cite{hartley};
    \item відношення просторової роздільності кадру та карти відоме з точністю до множника $\approx 2$ (див. п.~\ref{sec:scale});
    \item ортофотоплан не має географічної прив'язки, тому всі координати та похибки обчислюються в пікселях карти (\px{}), без переведення в метри.
\end{itemize}

Загальну схему обробки наведено на рисунку~\ref{fig:pipeline}. Система складається з офлайн-етапу підготовки карти, який виконується один раз, та онлайн-етапу обробки кадрів. Ключовим елементом онлайн-етапу є єдиний конвеєр зіставлення пари <<кадр --- фрагмент карти>>, який однаково застосовується в усіх режимах роботи: під час перевірки відомої пари, автоматичного пошуку фрагмента та трекінгу послідовності.

\begin{figure}[H]
\centering
\begin{tikzpicture}[
    node distance=5mm and 6mm,
    box/.style={draw, rounded corners=2pt, align=center, font=\small,
                minimum width=42mm, minimum height=9mm, inner sep=3pt},
    off/.style={box, fill=gray!12},
    on/.style={box, fill=blue!7},
    dec/.style={box, fill=orange!12},
    arr/.style={-{Stealth[length=2.2mm]}, thick}]
  % офлайн-етап
  \node[off] (map) {Ортофотоплан\\$6200\times4320$};
  \node[off, below=of map] (tiles) {Розбиття на фрагменти\\$512\times512$, перекриття 50\,\%};
  \node[off, below=of tiles] (cache) {Кеш ознак фрагментів\\(SuperPoint)};
  % онлайн-етап
  \node[on, right=16mm of map] (frame) {Кадр БпЛА\\$1920\times1080$};
  \node[on, below=of frame] (prep) {Нормалізація масштабу\\і компенсація курсу};
  \node[dec, below=of prep] (pool) {Пул фрагментів:\\вся карта / вікно $3\times3$};
  \node[on, below=of pool] (match) {Зіставлення\\LightGlue / LoFTR};
  \node[on, below=of match] (geom) {RANSAC + гомографія,\\перевірка правдоподібності};
  \node[dec, below=of geom] (rank) {Фільтр і ранжування\\фрагментів};
  \node[on, below=of rank] (coord) {Координати $(X,Y)$\\і контур огляду};
  \node[dec, below=of coord] (track) {Кінематична перевірка,\\траєкторія};
  \draw[arr] (map) -- (tiles);
  \draw[arr] (tiles) -- (cache);
  \draw[arr] (frame) -- (prep);
  \draw[arr] (prep) -- (pool);
  \draw[arr] (pool) -- (match);
  \draw[arr] (match) -- (geom);
  \draw[arr] (geom) -- (rank);
  \draw[arr] (rank) -- (coord);
  \draw[arr] (coord) -- (track);
  \draw[arr] (cache.south) |- (match.west);
  \draw[arr, dashed] (track.east) -- ++(6mm,0) |- node[pos=0.25, right, font=\footnotesize, align=left]{прогноз\\положення\\і курсу} (pool.east);
\end{tikzpicture}
\caption{Загальна схема роботи системи візуальної локалізації (сірим позначено офлайн-етап)}
\label{fig:pipeline}
\end{figure}

\subsection{Використані програмні засоби}

Систему реалізовано мовою Python~3.11. Використано такі бібліотеки:
\begin{itemize}[label=\textbf{•}]
    \item \textbf{PyTorch} 2.14 з підтримкою CUDA~12.6 \cite{pytorch} --- виконання нейромережевих моделей на GPU;
    \item \textbf{LightGlue} (офіційна реалізація авторів) --- моделі SuperPoint \cite{superpoint} і LightGlue \cite{lightglue};
    \item \textbf{Kornia} 0.8 \cite{kornia} --- реалізація LoFTR \cite{loftr} з попередньо навченими вагами \textit{outdoor};
    \item \textbf{OpenCV} 5.0 \cite{opencv} --- читання та масштабування зображень, оцінювання гомографії методами RANSAC \cite{ransac} і MAGSAC++ \cite{magsac}, геометричні перетворення;
    \item \textbf{NumPy} і \textbf{pandas} --- числові обчислення, робота з таблицями метаданих і результатів;
    \item \textbf{Matplotlib} --- побудова графіків;
    \item \textbf{pytest} --- автоматичне тестування модулів.
\end{itemize}

Нейромережі в роботі не навчалися: використано опубліковані попередньо навчені ваги моделей.

\subsection{Тестове середовище}

Розроблення та експерименти виконувалися на ноутбуці з такими характеристиками:
\begin{itemize}[label=\textbf{--}]
    \item \textbf{процесор:} AMD Ryzen 5 4600H;
    \item \textbf{оперативна пам'ять:} 8 ГБ;
    \item \textbf{відеокарта:} NVIDIA GeForce GTX 1650 (4 ГБ відеопам'яті);
    \item \textbf{CUDA:} 12.6 (збірка PyTorch \texttt{cu126});
    \item \textbf{операційна система:} Windows 11.
\end{itemize}

\subsection{Вхідні дані}

Для моделювання використано ортофотоплан району розміром $6200 \times 4320$ пікселів (формат RGBA з повністю непрозорим альфа-каналом) і відеозапис польоту БпЛА DJI (файл \code{DJI\_0331.MP4}): роздільність $1920 \times 1080$, частота 29,97~кадр/с, тривалість $\approx 94$~с, кодек HEVC, камера спрямована вниз.

Відео та ортофотоплан знято в різні пори року: ортофотоплан --- влітку, відео --- взимку. Через це рослинність, тіні та покриття ґрунту на кадрах і на карті суттєво відрізняються, що ускладнює зіставлення.

\subsection{Підготовка ортофотоплану}

\textbf{Розбиття на фрагменти з перекриттям.}\\
Нейромережеві методи зіставлення розраховані на зображення розміром порядку сотень пікселів, тому ортофотоплан розбивається на квадратні фрагменти (тайли) розміром $S = 512$ пікселів із часткою перекриття $o = 0{,}5$. Крок сітки визначається як
\begin{equation}
s = \lfloor S \cdot (1 - o) \rfloor = 256 .
\end{equation}
Координати лівих верхніх кутів фрагментів уздовж осі $x$ утворюють послідовність $0, s, 2s, \dots$; якщо регулярна сітка не доходить до правого краю карти, додається ще один крайовий фрагмент із початком у $W - S$ (аналогічно для осі $y$). Завдяки цьому всі фрагменти мають однаковий розмір і кожен піксель карти покрито. Кількість фрагментів:
\begin{equation}
N = \left( \left\lceil \frac{W - S}{s} \right\rceil + 1 \right) \cdot \left( \left\lceil \frac{H - S}{s} \right\rceil + 1 \right) = 24 \cdot 16 = 384 .
\end{equation}
Для кожного фрагмента в таблицю метаданих \code{tiles.csv} записуються ідентифікатор і координати $(X_0, Y_0)$ лівого верхнього кута в пікселях повної карти. Ці координати є основою для переходу від локальних координат фрагмента до глобальних координат карти.

Перекриття 50\,\% гарантує, що будь-яка ділянка карти, менша за половину фрагмента, повністю потрапляє принаймні в один фрагмент. Водночас, як показано в розділі~\ref{sec:experiments}, контур огляду камери на карті має розмір $\approx 518 \times 292$~\px{} і є ширшим за фрагмент, тому кадр, як правило, лише частково перекривається з кожним окремим фрагментом.

\textbf{Перевірка цілісності даних.}\\
Після розбиття автоматично перевіряється, що кожен фрагмент попіксельно збігається з відповідною ділянкою вихідної карти, усі фрагменти лежать у межах карти, їх кількість відповідає формулі, а таблиця метаданих не містить пропусків і дублікатів.

\subsection{Вилучення кадрів із відео}

З відео вибирається один кадр на секунду: кожен $n$-й кадр, де $n = \operatorname{round}(f_{\text{відео}} / f_{\text{ціль}}) = \operatorname{round}(29{,}97 / 1) = 30$. Часова мітка кадру обчислюється як $t = i / f_{\text{відео}}$, де $i$ --- номер кадру у відео. Кадри декодуються послідовно (без позиціювання за номером, яке ненадійне для кодеків H.264/HEVC) і зберігаються без втрат у форматі PNG. У результаті отримано 95 кадрів із кроком $\approx 1{,}001$~с.

\subsection{Нормалізація масштабу та єдина система координат}
\label{sec:scale}

Кадр і ортофотоплан мають різну просторову роздільність: один піксель кадру відповідає приблизно $s_f = 0{,}27$ пікселя карти (значення визначено калібруванням, див. п.~\ref{sec:exp-pair}). Якщо подати в метод зіставлення зображення з різним масштабом деталей, частина відповідностей втрачається. Тому перед зіставленням обидва зображення приводяться до спільної роздільності:
\begin{itemize}
    \item фрагмент карти масштабується з коефіцієнтом $k$ (у роботі $k = 1$, тобто фрагмент подається як $512 \times 512$);
    \item кадр масштабується з коефіцієнтом $s_f \cdot k$: кадр $1920 \times 1080$ перетворюється на $\approx 518 \times 292$ пікселі. Кадр використовується повністю, без обрізання.
\end{itemize}

Методи зіставлення повертають координати відповідностей у пікселях зменшених (<<робочих>>) зображень. Система переносить їх назад в оригінальні пікселі кадру та фрагмента з урахуванням узгодження центрів пікселів, яке використовується під час масштабування:
\begin{equation}
x_{\text{ориг}} = \frac{x_{\text{роб}} + 0{,}5}{s_x} - 0{,}5, \qquad
y_{\text{ориг}} = \frac{y_{\text{роб}} + 0{,}5}{s_y} - 0{,}5,
\label{eq:backmap}
\end{equation}
де $s_x, s_y$ --- фактичні коефіцієнти масштабування за осями (вони можуть дещо відрізнятися від заданого через округлення розмірів). Унаслідок цього гомографія $H$ оцінюється між оригінальними пікселями кадру та оригінальними пікселями фрагмента. Усі геометричні величини (поріг RANSAC, похибка репроєкції, координати центра й контуру кадру) виражаються в пікселях карти і не залежать ні від методу зіставлення, ні від робочого масштабу.

\subsection{Реалізація методів зіставлення}

\textbf{Спільний інтерфейс.}\\
Обидва методи реалізовано як класи зі спільним інтерфейсом, що складається з двох кроків:
\begin{enumerate}
    \item \code{extract(image)} --- підготовка одного зображення: для LightGlue --- виділення ключових точок і дескрипторів SuperPoint; для LoFTR, який не має окремого етапу детекції, --- лише формування тензора;
    \item \code{match\_features(f\textsubscript{0}, f\textsubscript{1})} --- встановлення відповідностей між двома підготовленими зображеннями.
\end{enumerate}
Таке розділення дає змогу обчислювати ознаки кадру один раз під час перебору багатьох фрагментів, а ознаки фрагментів карти --- заздалегідь. Решта системи працює лише з цим інтерфейсом і не залежить від конкретного методу.

\textbf{LightGlue + SuperPoint.}\\
SuperPoint виділяє до 2048 ключових точок на зображення; LightGlue використовується з параметрами авторів: поріг адаптивної зупинки 0,95, поріг відсікання точок 0,99, поріг прийняття відповідності 0,1. Під час аналізу прототипу курсової роботи виявлено, що функція виділення ознак SuperPoint за замовчуванням масштабує вхідне зображення до 1024~пікселів за довшою стороною. Унаслідок цього в курсовій роботі LightGlue фактично працював із зображеннями $1024 \times 1024$, а LoFTR --- з $512 \times 512$, тобто методи порівнювалися за різних умов. У розробленій системі це неявне масштабування вимкнено, і масштаб для обох методів задається однаково (п.~\ref{sec:scale}).

\textbf{LoFTR.}\\
Використано реалізацію бібліотеки Kornia з вагами \textit{outdoor}; зображення перетворюються у відтінки сірого. Грубий рівень LoFTR працює на сітці з кроком 8 пікселів, тому розміри вхідних зображень округлюються до кратних 8; відповідні фактичні коефіцієнти масштабу враховуються у формулі~(\ref{eq:backmap}).

\textbf{Кешування ознак фрагментів карти.}\\
Для LightGlue ознаки SuperPoint усіх 384 фрагментів обчислюються один раз і зберігаються на диску (координати точок, оцінки та дескриптори у форматі float32, без втрати точності): $523$--$1609$ точок на фрагмент, загальний обсяг $\approx 347$~МБ, час побудови $\approx 14$~с. Разом із кешем зберігаються параметри, з якими його побудовано (параметри SuperPoint, робочий масштаб, контрольні суми таблиці фрагментів і карти); кеш, що не відповідає поточним параметрам, система відмовляється використовувати. Для LoFTR кешування неможливе, оскільки його ознаки залежать від обох зображень пари.

\textbf{Вимірювання часу.}\\
Перед вимірюваннями виконується <<прогрівання>> моделей на зображеннях тих самих розмірів; для врахування асинхронних обчислень на GPU кожен вимірюваний блок обрамлюється синхронізацією CUDA. Час розділено на етапи: обробка кадру, обробка фрагмента, зіставлення, геометрична перевірка. Читання зображень із диска до вимірюваного часу не входить.

\subsection{Геометрична перевірка відповідностей}

\textbf{Оцінювання гомографії.}\\
За отриманими відповідностями $\{(\mathbf{p}_i, \mathbf{q}_i)\}$ (точка кадру --- точка фрагмента) оцінюється гомографія $H$ (див. розділ~1). Використовується алгоритм RANSAC \cite{ransac} із порогом $\tau = 3$~\px{}, максимальною кількістю ітерацій 5000 і довірчою ймовірністю 0,999; для відтворюваності генератор випадкових чисел ініціалізується фіксованим значенням. Як альтернативний оцінювач реалізовано MAGSAC++ \cite{magsac}, який замість жорсткого порогу зважує відповідності за їхньою похибкою; порівняння оцінювачів наведено в розділі~\ref{sec:experiments}.

Відповідність вважається \emph{геометрично узгодженою}, якщо пряма похибка репроєкції не перевищує порогу:
\begin{equation}
e_i = \left\| \pi(H \tilde{\mathbf{p}}_i) - \mathbf{q}_i \right\|_2 \le \tau,
\end{equation}
де $\tilde{\mathbf{p}}_i$ --- однорідні координати точки, $\pi([u, v, w]^T) = (u/w,\, v/w)$ --- перехід від однорідних координат до декартових. Якщо кількість відповідностей менша за мінімально необхідну для успіху (15), RANSAC не запускається: на наборі, що складається лише з хибних пар, \code{cv2.RANSAC} відпрацьовує всі 5000 ітерацій ($\approx 0{,}1$~с) без шансу на успішний результат.

\textbf{Перевірка правдоподібності гомографії.}\\
RANSAC може знайти формально узгоджену, але фізично неможливу модель (наприклад, на повторюваних структурах). Тому для знайденої $H$ обчислюється якобіан $J$ відображення в центрі кадру $\mathbf{c}$ --- локальне афінне наближення гомографії --- і з нього визначаються:
\begin{align}
\sigma &= \sqrt{|\det J|} && \text{--- масштаб (пікселів карти на піксель кадру)}, \\
\theta &= \operatorname{atan2}(J_{21} - J_{12},\; J_{11} + J_{22}) && \text{--- кут повороту кадру відносно карти}, \\
\kappa &= \sigma_{\max}(J) / \sigma_{\min}(J) && \text{--- анізотропія (відношення сингулярних чисел)}.
\end{align}
Гомографія визнається правдоподібною, якщо виконуються всі умови:
\begin{enumerate}
    \item третя однорідна координата має однаковий знак в усіх кутах кадру (рамка кадру не перетинає <<лінію нескінченності>>);
    \item $\det J > 0$ --- немає дзеркального відображення;
    \item проєкція рамки кадру є опуклим чотирикутником;
    \item $\kappa \le 1{,}5$ --- для надирної камери масштаб за осями майже однаковий;
    \item $s_f / 2 \le \sigma \le 2 s_f$ --- масштаб узгоджується з очікуваним.
\end{enumerate}
Кут $\theta$ надалі використовується як оцінка курсу БпЛА відносно карти, а масштаб $\sigma$ --- як непрямий показник висоти польоту.

\textbf{Критерій успішного зіставлення пари.}\\
Зіставлення пари <<кадр --- фрагмент>> вважається успішним, якщо гомографію знайдено, кількість геометрично узгоджених відповідностей $N_{\text{in}} \ge 15$ і гомографія правдоподібна. Інакше система фіксує конкретну причину невдачі.

\subsection{Автоматичний пошук фрагмента карти}
\label{sec:search}

Для кадру, положення якого невідоме, конвеєр зіставлення виконується з кожним фрагментом пулу (усієї карти, заданої прямокутної області або прорідженої сітки фрагментів). Ознаки кадру обчислюються один раз. Вибір найкращого фрагмента виконується у два кроки.

\textbf{Крок 1: жорсткий фільтр.}
Фрагмент-кандидат відкидається, якщо не виконується хоча б одна умова:
\begin{enumerate}
    \item гомографію знайдено і вона правдоподібна;
    \item $N_{\text{in}} \ge 15$;
    \item частка геометрично узгоджених відповідностей $r = N_{\text{in}} / N_{\text{all}} \ge 0{,}15$;
    \item середня похибка репроєкції $\bar{e} \le 3{,}5$~\px{};
    \item проєкція центра кадру $(x_{\text{лок}}, y_{\text{лок}})$ потрапляє в межі фрагмента з допуском $m = 30$~\px{}: $-m \le x_{\text{лок}} \le S - 1 + m$ (аналогічно для $y$), тобто фрагмент є основним носієм центра кадру;
    \item проєкція центра кадру лежить у межах карти.
\end{enumerate}
Умова~4 за порогу RANSAC $\tau = 3$~\px{} виконується завжди і є запобіжником на випадок зміни порогу.

\textbf{Крок 2: бал ранжування.}
Серед кандидатів, що пройшли фільтр, обирається фрагмент із найбільшим значенням
\begin{equation}
\operatorname{Score} = N_{\text{in}} \cdot \frac{N_{\text{in}}}{N_{\text{all}}} \cdot \exp\!\left(-\frac{\bar{e}}{\sigma_r}\right), \qquad \sigma_r = 2~\text{\px{}}.
\label{eq:score}
\end{equation}
Бал зростає з кількістю та часткою геометрично узгоджених відповідностей і плавно зменшується з похибкою репроєкції. На відміну від відношення $N_{\text{in}} / \bar{e}$, він не стає необмежено великим за малої похибки на кількох випадкових точках.

\textbf{Обчислення координат на карті.}\\
Центр кадру $\mathbf{c} = ((W_f - 1)/2,\, (H_f - 1)/2)$ (у конвенції центрів пікселів) переноситься гомографією обраного фрагмента з явною дегомогенізацією:
\begin{equation}
\tilde{\mathbf{p}} = H \begin{bmatrix} c_x \\ c_y \\ 1 \end{bmatrix}, \qquad
x_{\text{лок}} = \frac{\tilde{p}_1}{\tilde{p}_3}, \quad y_{\text{лок}} = \frac{\tilde{p}_2}{\tilde{p}_3},
\end{equation}
після чого обчислюються глобальні координати на карті:
\begin{equation}
X = X_0 + x_{\text{лок}}, \qquad Y = Y_0 + y_{\text{лок}},
\end{equation}
де $(X_0, Y_0)$ --- координати кута фрагмента з таблиці метаданих. Аналогічно переносяться чотири кути кадру, що дає контур огляду камери на карті. Точки, для яких $|\tilde{p}_3| \approx 0$, вважаються невизначеними й відкидаються.

\subsection{Локалізація послідовності кадрів}
\label{sec:tracking}

Повний перебір 384 фрагментів для кожного кадру є надлишковим: за 1~с апарат зміщується лише на кілька десятків пікселів карти. Тому для послідовності кадрів реалізовано трекер зі станами \textsc{lost} та \textsc{tracking}.

\textbf{Стан \textsc{lost}} (перший кадр або втрата супроводу): виконується глобальний пошук (п.~\ref{sec:search}). Якщо переможець має бал $\operatorname{Score} \ge 5$, положення приймається (статус \textsc{initialized}) і трекер переходить у стан \textsc{tracking}.

\textbf{Стан \textsc{tracking}}: пошук виконується лише серед фрагментів вікна $3 \times 3$ сітки (не більше 9 фрагментів) навколо базового фрагмента. Базовим обирається фрагмент, центр якого найближчий до \emph{прогнозованого} положення. Прогноз будується за двома останніми прийнятими положеннями $(t_{k-2}, P_{k-2})$, $(t_{k-1}, P_{k-1})$ у припущенні сталої швидкості:
\begin{equation}
\hat{P}_k = P_{k-1} + \mathbf{v} \cdot (t_k - t_{k-1}), \qquad
\mathbf{v} = \frac{P_{k-1} - P_{k-2}}{t_{k-1} - t_{k-2}},
\label{eq:extrapolation}
\end{equation}
де модуль швидкості обмежується значенням $V_{\max}$. За сталого інтервалу між кадрами це збігається з $\hat{P}_k = P_{k-1} + (P_{k-1} - P_{k-2})$.

\textbf{Кінематична перевірка.}\\
Нехай $P_{\text{ост}}$ --- останнє \emph{прийняте} положення, отримане в момент $t_{\text{ост}}$. Нова оцінка $P_k$ приймається, якщо
\begin{equation}
\| P_k - P_{\text{ост}} \|_2 \le D_{\max} = V_{\max} \cdot (t_k - t_{\text{ост}}) + D_{\text{запас}},
\end{equation}
де $V_{\max} = 120$~\px{}/с, $D_{\text{запас}} = 30$~\px{}; за частоти 1~кадр/с $D_{\max} = 150$~\px{}. Порівняння саме з останнім прийнятим положенням означає, що після кількох пропущених кадрів допустимий радіус зростає пропорційно часу. Серед кандидатів, що пройшли фільтр п.~\ref{sec:search}, обирається найкращий за балом, який проходить і кінематичну перевірку. Якщо такого кандидата немає, положення кадру екстраполюється за формулою~(\ref{eq:extrapolation}) (статус \textsc{interpolated}), а відкинута оцінка зберігається для аналізу. Після трьох поспіль кадрів без прийнятого положення трекер переходить у стан \textsc{lost}.

\textbf{Компенсація курсу.}\\
SuperPoint/LightGlue і LoFTR не інваріантні до великих поворотів зображення (кількісне підтвердження наведено в п.~\ref{sec:exp-robust}), а під час розвороту БпЛА кадр повертається відносно карти. Кут <<кадр $\to$ карта>> відомий із гомографій прийнятих кадрів ($\theta$, п.~\ref{sec:search}), тому перед зіставленням кадр повертається на прогнозований кут $\hat{\theta}_k$, і метод зіставлення отримує кадр, орієнтований так само, як карта. Прогноз кута будується аналогічно до~(\ref{eq:extrapolation}) з урахуванням переходу через $\pm 180^\circ$, а кутова швидкість обмежується значенням $45^\circ$/с. Полотно під час повороту розширюється, а відповідності, що потрапили на фон поза кадром, відкидаються. Перетворення <<оригінальний кадр $\to$ робоче зображення>> є композицією масштабування та повороту $T = R(\hat{\theta}_k) \cdot S$, і координати відповідностей переносяться назад точним оберненим перетворенням $T^{-1}$. Тому гомографія, як і раніше, оцінюється в оригінальних пікселях кадру.

Якщо з прогнозованим курсом фрагмент не знайдено, у стані \textsc{tracking} послідовно перевіряються гіпотези $\hat{\theta}_k \pm 45^\circ$ і $\hat{\theta}_k \pm 90^\circ$. У стані \textsc{lost}, де повний перебір карти дорогий, кожен наступний глобальний пошук використовує чергову гіпотезу з циклу $0^\circ, 90^\circ, 180^\circ, -90^\circ$ відносно останнього відомого курсу.

Узагальнено обробка одного кадру в стані \textsc{tracking} складається з таких кроків:
\begin{enumerate}
    \item обчислити прогноз положення $\hat{P}_k$ і курсу $\hat{\theta}_k$;
    \item сформувати пул із не більше ніж 9 фрагментів навколо фрагмента, найближчого до $\hat{P}_k$;
    \item для кожної гіпотези курсу, починаючи з $\hat{\theta}_k$:
    \begin{enumerate}
        \item повернути та масштабувати кадр, обчислити його ознаки;
        \item зіставити кадр із фрагментами пулу, виконати геометричну перевірку, фільтр і ранжування;
        \item якщо найкращий кандидат проходить кінематичну перевірку --- прийняти положення (статус \textsc{tracking}) і завершити обробку кадру;
    \end{enumerate}
    \item якщо жодна гіпотеза не дала прийнятного положення --- екстраполювати положення (статус \textsc{interpolated}); після трьох таких кадрів поспіль перейти в стан \textsc{lost}.
\end{enumerate}

\subsection{Модульна структура та відтворюваність}

Програмну систему організовано як пакет із незалежними модулями (табл.~\ref{tab:modules}) і набір сценаріїв командного рядка для кожного етапу.

\begin{table}[H]
\centering
\caption{Основні модулі програмної системи}
\label{tab:modules}
\small
\begin{tabular}{p{0.36\linewidth} p{0.56\linewidth}}
\toprule
\textbf{Модуль} & \textbf{Призначення} \\
\midrule
\code{preprocessing/tiles.py}, \code{frames.py} & розбиття карти, вилучення кадрів, метадані \\
\code{preprocessing/images.py} & читання/запис, масштабування, поворот, перенесення координат, спотворення для експериментів \\
\code{preprocessing/cache.py} & дисковий кеш ознак фрагментів \\
\code{matching/lightglue.py}, \code{loftr.py} & реалізації методів зі спільним інтерфейсом \\
\code{matching/pipeline.py} & єдиний конвеєр <<кадр + фрагмент $\to$ відповідності $\to$ RANSAC $\to$ метрики>> \\
\code{geometry/homography.py} & RANSAC/MAGSAC++, похибка репроєкції, перевірка правдоподібності \\
\code{evaluation/scoring.py} & фільтр і бал ранжування фрагментів \\
\code{localization/coordinates.py} & перенесення точок на карту, контур огляду \\
\code{localization/search.py}, \code{tracker.py} & пошук фрагмента, трекер послідовності \\
\code{evaluation/metrics.py}, \code{ground\_truth.py} & метрики, еталонні координати \\
\code{visualization/} & рисунки відповідностей, положення та траєкторій \\
\bottomrule
\end{tabular}
\end{table}

Для відтворюваності результатів кожен запуск записує, разом із результатами, файл метаданих з усіма параметрами конвеєра й методу зіставлення, ідентифікатором версії коду та відомостями про обчислювальне середовище. Генератори випадкових чисел RANSAC і спотворень ініціалізуються фіксованими значеннями. Коректність перенесення координат, оцінювання гомографії, критеріїв фільтрації та логіки трекера перевіряється набором із 72 автоматичних тестів на синтетичних даних, для яких точний результат відомий заздалегідь; тести не потребують нейромереж і виконуються за кілька секунд.

\subsection{Візуалізація результатів}

Система формує такі графічні матеріали:
\begin{itemize}
    \item відповідності між кадром і фрагментом карти: геометрично узгоджені --- зеленим, решта --- червоним;
    \item контур огляду камери і центр кадру на ортофотоплані разом із сіткою перевірених фрагментів і обраним фрагментом;
    \item траєкторію польоту на карті з позначенням статусів кадрів і напрямку руху;
    \item графіки похибок, часу обробки та результатів експериментів (розділ~\ref{sec:experiments}).
\end{itemize}

% Потрібно в преамбулі: \usetikzlibrary{arrows.meta,positioning,babel}
% =============================================================================
\providecommand{\px}{пкс}               % пікселі карти
\providecommand{\code}[1]{\texttt{#1}}  % імена файлів, модулів і функцій
\newpage
\section{Програмне забезпечення та програмна реалізація системи}
\label{sec:implementation}

\subsection{Постановка задачі та загальна архітектура системи}

Розроблена система розв'язує задачу абсолютної візуальної локалізації БпЛА: за кадром із бортової камери та ортофотопланом місцевості визначити, над якою точкою карти перебуває апарат, а за послідовністю кадрів --- відновити траєкторію польоту. Інформація про те, якому фрагменту карти відповідає кадр, системі не надається: відповідну ділянку ортофотоплану система знаходить автоматично.

Вхідними даними є:
\begin{itemize}
    \item ортофотоплан місцевості розміром $W \times H$ пікселів;
    \item відеозапис польоту БпЛА або послідовність кадрів із часовими мітками.
\end{itemize}

Вихідними даними є:
\begin{itemize}
    \item для кожного кадру --- координати $(X, Y)$ проєкції центра кадру на ортофотоплан і контур огляду камери (чотирикутник на карті);
    \item для послідовності кадрів --- траєкторія польоту з позначкою статусу кожного кадру та показниками якості зіставлення.
\end{itemize}

Робота системи ґрунтується на таких припущеннях:
\begin{itemize}
    \item камера спрямована приблизно вертикально вниз (надирна зйомка);
    \item сцену в межах кадру можна наближено вважати плоскою, тому перетворення <<кадр $\to$ карта>> описується гомографією \cite{hartley};
    \item відношення просторової роздільності кадру та карти відоме з точністю до множника $\approx 2$ (див. п.~\ref{sec:scale});
    \item ортофотоплан не має географічної прив'язки, тому всі координати та похибки обчислюються в пікселях карти (\px{}), без переведення в метри.
\end{itemize}

Загальну схему обробки наведено на рисунку~\ref{fig:pipeline}. Система складається з офлайн-етапу підготовки карти, який виконується один раз, та онлайн-етапу обробки кадрів. Ключовим елементом онлайн-етапу є єдиний конвеєр зіставлення пари <<кадр --- фрагмент карти>>, який однаково застосовується в усіх режимах роботи: під час перевірки відомої пари, автоматичного пошуку фрагмента та трекінгу послідовності.

\begin{figure}[H]
\centering
\begin{tikzpicture}[
    node distance=5mm and 6mm,
    box/.style={draw, rounded corners=2pt, align=center, font=\small,
                minimum width=42mm, minimum height=9mm, inner sep=3pt},
    off/.style={box, fill=gray!12},
    on/.style={box, fill=blue!7},
    dec/.style={box, fill=orange!12},
    arr/.style={-{Stealth[length=2.2mm]}, thick}]
  % офлайн-етап
  \node[off] (map) {Ортофотоплан\\$6200\times4320$};
  \node[off, below=of map] (tiles) {Розбиття на фрагменти\\$512\times512$, перекриття 50\,\%};
  \node[off, below=of tiles] (cache) {Кеш ознак фрагментів\\(SuperPoint)};
  % онлайн-етап
  \node[on, right=16mm of map] (frame) {Кадр БпЛА\\$1920\times1080$};
  \node[on, below=of frame] (prep) {Нормалізація масштабу\\і компенсація курсу};
  \node[dec, below=of prep] (pool) {Пул фрагментів:\\вся карта / вікно $3\times3$};
  \node[on, below=of pool] (match) {Зіставлення\\LightGlue / LoFTR};
  \node[on, below=of match] (geom) {RANSAC + гомографія,\\перевірка правдоподібності};
  \node[dec, below=of geom] (rank) {Фільтр і ранжування\\фрагментів};
  \node[on, below=of rank] (coord) {Координати $(X,Y)$\\і контур огляду};
  \node[dec, below=of coord] (track) {Кінематична перевірка,\\траєкторія};
  \draw[arr] (map) -- (tiles);
  \draw[arr] (tiles) -- (cache);
  \draw[arr] (frame) -- (prep);
  \draw[arr] (prep) -- (pool);
  \draw[arr] (pool) -- (match);
  \draw[arr] (match) -- (geom);
  \draw[arr] (geom) -- (rank);
  \draw[arr] (rank) -- (coord);
  \draw[arr] (coord) -- (track);
  \draw[arr] (cache.south) |- (match.west);
  \draw[arr, dashed] (track.east) -- ++(6mm,0) |- node[pos=0.25, right, font=\footnotesize, align=left]{прогноз\\положення\\і курсу} (pool.east);
\end{tikzpicture}
\caption{Загальна схема роботи системи візуальної локалізації (сірим позначено офлайн-етап)}
\label{fig:pipeline}
\end{figure}

\subsection{Використані програмні засоби}

Систему реалізовано мовою Python~3.11. Використано такі бібліотеки:
\begin{itemize}[label=\textbf{•}]
    \item \textbf{PyTorch} 2.14 з підтримкою CUDA~12.6 \cite{pytorch} --- виконання нейромережевих моделей на GPU;
    \item \textbf{LightGlue} (офіційна реалізація авторів) --- моделі SuperPoint \cite{superpoint} і LightGlue \cite{lightglue};
    \item \textbf{Kornia} 0.8 \cite{kornia} --- реалізація LoFTR \cite{loftr} з попередньо навченими вагами \textit{outdoor};
    \item \textbf{OpenCV} 5.0 \cite{opencv} --- читання та масштабування зображень, оцінювання гомографії методами RANSAC \cite{ransac} і MAGSAC++ \cite{magsac}, геометричні перетворення;
    \item \textbf{NumPy} і \textbf{pandas} --- числові обчислення, робота з таблицями метаданих і результатів;
    \item \textbf{Matplotlib} --- побудова графіків;
    \item \textbf{pytest} --- автоматичне тестування модулів.
\end{itemize}

Нейромережі в роботі не навчалися: використано опубліковані попередньо навчені ваги моделей.

\subsection{Тестове середовище}

Розроблення та експерименти виконувалися на ноутбуці з такими характеристиками:
\begin{itemize}[label=\textbf{--}]
    \item \textbf{процесор:} AMD Ryzen 5 4600H;
    \item \textbf{оперативна пам'ять:} 8 ГБ;
    \item \textbf{відеокарта:} NVIDIA GeForce GTX 1650 (4 ГБ відеопам'яті);
    \item \textbf{CUDA:} 12.6 (збірка PyTorch \texttt{cu126});
    \item \textbf{операційна система:} Windows 11.
\end{itemize}

\subsection{Вхідні дані}

Для моделювання використано ортофотоплан району розміром $6200 \times 4320$ пікселів (формат RGBA з повністю непрозорим альфа-каналом) і відеозапис польоту БпЛА DJI (файл \code{DJI\_0331.MP4}): роздільність $1920 \times 1080$, частота 29,97~кадр/с, тривалість $\approx 94$~с, кодек HEVC, камера спрямована вниз.

Відео та ортофотоплан знято в різні пори року: ортофотоплан --- влітку, відео --- взимку. Через це рослинність, тіні та покриття ґрунту на кадрах і на карті суттєво відрізняються, що ускладнює зіставлення.

\subsection{Підготовка ортофотоплану}

\textbf{Розбиття на фрагменти з перекриттям.}\\
Нейромережеві методи зіставлення розраховані на зображення розміром порядку сотень пікселів, тому ортофотоплан розбивається на квадратні фрагменти (тайли) розміром $S = 512$ пікселів із часткою перекриття $o = 0{,}5$. Крок сітки визначається як
\begin{equation}
s = \lfloor S \cdot (1 - o) \rfloor = 256 .
\end{equation}
Координати лівих верхніх кутів фрагментів уздовж осі $x$ утворюють послідовність $0, s, 2s, \dots$; якщо регулярна сітка не доходить до правого краю карти, додається ще один крайовий фрагмент із початком у $W - S$ (аналогічно для осі $y$). Завдяки цьому всі фрагменти мають однаковий розмір і кожен піксель карти покрито. Кількість фрагментів:
\begin{equation}
N = \left( \left\lceil \frac{W - S}{s} \right\rceil + 1 \right) \cdot \left( \left\lceil \frac{H - S}{s} \right\rceil + 1 \right) = 24 \cdot 16 = 384 .
\end{equation}
Для кожного фрагмента в таблицю метаданих \code{tiles.csv} записуються ідентифікатор і координати $(X_0, Y_0)$ лівого верхнього кута в пікселях повної карти. Ці координати є основою для переходу від локальних координат фрагмента до глобальних координат карти.

Перекриття 50\,\% гарантує, що будь-яка ділянка карти, менша за половину фрагмента, повністю потрапляє принаймні в один фрагмент. Водночас, як показано в розділі~\ref{sec:experiments}, контур огляду камери на карті має розмір $\approx 518 \times 292$~\px{} і є ширшим за фрагмент, тому кадр, як правило, лише частково перекривається з кожним окремим фрагментом.

\textbf{Перевірка цілісності даних.}\\
Після розбиття автоматично перевіряється, що кожен фрагмент попіксельно збігається з відповідною ділянкою вихідної карти, усі фрагменти лежать у межах карти, їх кількість відповідає формулі, а таблиця метаданих не містить пропусків і дублікатів.

\subsection{Вилучення кадрів із відео}

З відео вибирається один кадр на секунду: кожен $n$-й кадр, де $n = \operatorname{round}(f_{\text{відео}} / f_{\text{ціль}}) = \operatorname{round}(29{,}97 / 1) = 30$. Часова мітка кадру обчислюється як $t = i / f_{\text{відео}}$, де $i$ --- номер кадру у відео. Кадри декодуються послідовно (без позиціювання за номером, яке ненадійне для кодеків H.264/HEVC) і зберігаються без втрат у форматі PNG. У результаті отримано 95 кадрів із кроком $\approx 1{,}001$~с.

\subsection{Нормалізація масштабу та єдина система координат}
\label{sec:scale}

Кадр і ортофотоплан мають різну просторову роздільність: один піксель кадру відповідає приблизно $s_f = 0{,}27$ пікселя карти (значення визначено калібруванням, див. п.~\ref{sec:exp-pair}). Якщо подати в метод зіставлення зображення з різним масштабом деталей, частина відповідностей втрачається. Тому перед зіставленням обидва зображення приводяться до спільної роздільності:
\begin{itemize}
    \item фрагмент карти масштабується з коефіцієнтом $k$ (у роботі $k = 1$, тобто фрагмент подається як $512 \times 512$);
    \item кадр масштабується з коефіцієнтом $s_f \cdot k$: кадр $1920 \times 1080$ перетворюється на $\approx 518 \times 292$ пікселі. Кадр використовується повністю, без обрізання.
\end{itemize}

Методи зіставлення повертають координати відповідностей у пікселях зменшених (<<робочих>>) зображень. Система переносить їх назад в оригінальні пікселі кадру та фрагмента з урахуванням узгодження центрів пікселів, яке використовується під час масштабування:
\begin{equation}
x_{\text{ориг}} = \frac{x_{\text{роб}} + 0{,}5}{s_x} - 0{,}5, \qquad
y_{\text{ориг}} = \frac{y_{\text{роб}} + 0{,}5}{s_y} - 0{,}5,
\label{eq:backmap}
\end{equation}
де $s_x, s_y$ --- фактичні коефіцієнти масштабування за осями (вони можуть дещо відрізнятися від заданого через округлення розмірів). Унаслідок цього гомографія $H$ оцінюється між оригінальними пікселями кадру та оригінальними пікселями фрагмента. Усі геометричні величини (поріг RANSAC, похибка репроєкції, координати центра й контуру кадру) виражаються в пікселях карти і не залежать ні від методу зіставлення, ні від робочого масштабу.

\subsection{Реалізація методів зіставлення}

\textbf{Спільний інтерфейс.}\\
Обидва методи реалізовано як класи зі спільним інтерфейсом, що складається з двох кроків:
\begin{enumerate}
    \item \code{extract(image)} --- підготовка одного зображення: для LightGlue --- виділення ключових точок і дескрипторів SuperPoint; для LoFTR, який не має окремого етапу детекції, --- лише формування тензора;
    \item \code{match\_features(f\textsubscript{0}, f\textsubscript{1})} --- встановлення відповідностей між двома підготовленими зображеннями.
\end{enumerate}
Таке розділення дає змогу обчислювати ознаки кадру один раз під час перебору багатьох фрагментів, а ознаки фрагментів карти --- заздалегідь. Решта системи працює лише з цим інтерфейсом і не залежить від конкретного методу.

\textbf{LightGlue + SuperPoint.}\\
SuperPoint виділяє до 2048 ключових точок на зображення; LightGlue використовується з параметрами авторів: поріг адаптивної зупинки 0,95, поріг відсікання точок 0,99, поріг прийняття відповідності 0,1. Під час аналізу прототипу курсової роботи виявлено, що функція виділення ознак SuperPoint за замовчуванням масштабує вхідне зображення до 1024~пікселів за довшою стороною. Унаслідок цього в курсовій роботі LightGlue фактично працював із зображеннями $1024 \times 1024$, а LoFTR --- з $512 \times 512$, тобто методи порівнювалися за різних умов. У розробленій системі це неявне масштабування вимкнено, і масштаб для обох методів задається однаково (п.~\ref{sec:scale}).

\textbf{LoFTR.}\\
Використано реалізацію бібліотеки Kornia з вагами \textit{outdoor}; зображення перетворюються у відтінки сірого. Грубий рівень LoFTR працює на сітці з кроком 8 пікселів, тому розміри вхідних зображень округлюються до кратних 8; відповідні фактичні коефіцієнти масштабу враховуються у формулі~(\ref{eq:backmap}).

\textbf{Кешування ознак фрагментів карти.}\\
Для LightGlue ознаки SuperPoint усіх 384 фрагментів обчислюються один раз і зберігаються на диску (координати точок, оцінки та дескриптори у форматі float32, без втрати точності): $523$--$1609$ точок на фрагмент, загальний обсяг $\approx 347$~МБ, час побудови $\approx 14$~с. Разом із кешем зберігаються параметри, з якими його побудовано (параметри SuperPoint, робочий масштаб, контрольні суми таблиці фрагментів і карти); кеш, що не відповідає поточним параметрам, система відмовляється використовувати. Для LoFTR кешування неможливе, оскільки його ознаки залежать від обох зображень пари.

\textbf{Вимірювання часу.}\\
Перед вимірюваннями виконується <<прогрівання>> моделей на зображеннях тих самих розмірів; для врахування асинхронних обчислень на GPU кожен вимірюваний блок обрамлюється синхронізацією CUDA. Час розділено на етапи: обробка кадру, обробка фрагмента, зіставлення, геометрична перевірка. Читання зображень із диска до вимірюваного часу не входить.

\subsection{Геометрична перевірка відповідностей}

\textbf{Оцінювання гомографії.}\\
За отриманими відповідностями $\{(\mathbf{p}_i, \mathbf{q}_i)\}$ (точка кадру --- точка фрагмента) оцінюється гомографія $H$ (див. розділ~1). Використовується алгоритм RANSAC \cite{ransac} із порогом $\tau = 3$~\px{}, максимальною кількістю ітерацій 5000 і довірчою ймовірністю 0,999; для відтворюваності генератор випадкових чисел ініціалізується фіксованим значенням. Як альтернативний оцінювач реалізовано MAGSAC++ \cite{magsac}, який замість жорсткого порогу зважує відповідності за їхньою похибкою; порівняння оцінювачів наведено в розділі~\ref{sec:experiments}.

Відповідність вважається \emph{геометрично узгодженою}, якщо пряма похибка репроєкції не перевищує порогу:
\begin{equation}
e_i = \left\| \pi(H \tilde{\mathbf{p}}_i) - \mathbf{q}_i \right\|_2 \le \tau,
\end{equation}
де $\tilde{\mathbf{p}}_i$ --- однорідні координати точки, $\pi([u, v, w]^T) = (u/w,\, v/w)$ --- перехід від однорідних координат до декартових. Якщо кількість відповідностей менша за мінімально необхідну для успіху (15), RANSAC не запускається: на наборі, що складається лише з хибних пар, \code{cv2.RANSAC} відпрацьовує всі 5000 ітерацій ($\approx 0{,}1$~с) без шансу на успішний результат.

\textbf{Перевірка правдоподібності гомографії.}\\
RANSAC може знайти формально узгоджену, але фізично неможливу модель (наприклад, на повторюваних структурах). Тому для знайденої $H$ обчислюється якобіан $J$ відображення в центрі кадру $\mathbf{c}$ --- локальне афінне наближення гомографії --- і з нього визначаються:
\begin{align}
\sigma &= \sqrt{|\det J|} && \text{--- масштаб (пікселів карти на піксель кадру)}, \\
\theta &= \operatorname{atan2}(J_{21} - J_{12},\; J_{11} + J_{22}) && \text{--- кут повороту кадру відносно карти}, \\
\kappa &= \sigma_{\max}(J) / \sigma_{\min}(J) && \text{--- анізотропія (відношення сингулярних чисел)}.
\end{align}
Гомографія визнається правдоподібною, якщо виконуються всі умови:
\begin{enumerate}
    \item третя однорідна координата має однаковий знак в усіх кутах кадру (рамка кадру не перетинає <<лінію нескінченності>>);
    \item $\det J > 0$ --- немає дзеркального відображення;
    \item проєкція рамки кадру є опуклим чотирикутником;
    \item $\kappa \le 1{,}5$ --- для надирної камери масштаб за осями майже однаковий;
    \item $s_f / 2 \le \sigma \le 2 s_f$ --- масштаб узгоджується з очікуваним.
\end{enumerate}
Кут $\theta$ надалі використовується як оцінка курсу БпЛА відносно карти, а масштаб $\sigma$ --- як непрямий показник висоти польоту.

\textbf{Критерій успішного зіставлення пари.}\\
Зіставлення пари <<кадр --- фрагмент>> вважається успішним, якщо гомографію знайдено, кількість геометрично узгоджених відповідностей $N_{\text{in}} \ge 15$ і гомографія правдоподібна. Інакше система фіксує конкретну причину невдачі.

\subsection{Автоматичний пошук фрагмента карти}
\label{sec:search}

Для кадру, положення якого невідоме, конвеєр зіставлення виконується з кожним фрагментом пулу (усієї карти, заданої прямокутної області або прорідженої сітки фрагментів). Ознаки кадру обчислюються один раз. Вибір найкращого фрагмента виконується у два кроки.

\textbf{Крок 1: жорсткий фільтр.}
Фрагмент-кандидат відкидається, якщо не виконується хоча б одна умова:
\begin{enumerate}
    \item гомографію знайдено і вона правдоподібна;
    \item $N_{\text{in}} \ge 15$;
    \item частка геометрично узгоджених відповідностей $r = N_{\text{in}} / N_{\text{all}} \ge 0{,}15$;
    \item середня похибка репроєкції $\bar{e} \le 3{,}5$~\px{};
    \item проєкція центра кадру $(x_{\text{лок}}, y_{\text{лок}})$ потрапляє в межі фрагмента з допуском $m = 30$~\px{}: $-m \le x_{\text{лок}} \le S - 1 + m$ (аналогічно для $y$), тобто фрагмент є основним носієм центра кадру;
    \item проєкція центра кадру лежить у межах карти.
\end{enumerate}
Умова~4 за порогу RANSAC $\tau = 3$~\px{} виконується завжди і є запобіжником на випадок зміни порогу.

\textbf{Крок 2: бал ранжування.}
Серед кандидатів, що пройшли фільтр, обирається фрагмент із найбільшим значенням
\begin{equation}
\operatorname{Score} = N_{\text{in}} \cdot \frac{N_{\text{in}}}{N_{\text{all}}} \cdot \exp\!\left(-\frac{\bar{e}}{\sigma_r}\right), \qquad \sigma_r = 2~\text{\px{}}.
\label{eq:score}
\end{equation}
Бал зростає з кількістю та часткою геометрично узгоджених відповідностей і плавно зменшується з похибкою репроєкції. На відміну від відношення $N_{\text{in}} / \bar{e}$, він не стає необмежено великим за малої похибки на кількох випадкових точках.

\textbf{Обчислення координат на карті.}\\
Центр кадру $\mathbf{c} = ((W_f - 1)/2,\, (H_f - 1)/2)$ (у конвенції центрів пікселів) переноситься гомографією обраного фрагмента з явною дегомогенізацією:
\begin{equation}
\tilde{\mathbf{p}} = H \begin{bmatrix} c_x \\ c_y \\ 1 \end{bmatrix}, \qquad
x_{\text{лок}} = \frac{\tilde{p}_1}{\tilde{p}_3}, \quad y_{\text{лок}} = \frac{\tilde{p}_2}{\tilde{p}_3},
\end{equation}
після чого обчислюються глобальні координати на карті:
\begin{equation}
X = X_0 + x_{\text{лок}}, \qquad Y = Y_0 + y_{\text{лок}},
\end{equation}
де $(X_0, Y_0)$ --- координати кута фрагмента з таблиці метаданих. Аналогічно переносяться чотири кути кадру, що дає контур огляду камери на карті. Точки, для яких $|\tilde{p}_3| \approx 0$, вважаються невизначеними й відкидаються.

\subsection{Локалізація послідовності кадрів}
\label{sec:tracking}

Повний перебір 384 фрагментів для кожного кадру є надлишковим: за 1~с апарат зміщується лише на кілька десятків пікселів карти. Тому для послідовності кадрів реалізовано трекер зі станами \textsc{lost} та \textsc{tracking}.

\textbf{Стан \textsc{lost}} (перший кадр або втрата супроводу): виконується глобальний пошук (п.~\ref{sec:search}). Якщо переможець має бал $\operatorname{Score} \ge 5$, положення приймається (статус \textsc{initialized}) і трекер переходить у стан \textsc{tracking}.

\textbf{Стан \textsc{tracking}}: пошук виконується лише серед фрагментів вікна $3 \times 3$ сітки (не більше 9 фрагментів) навколо базового фрагмента. Базовим обирається фрагмент, центр якого найближчий до \emph{прогнозованого} положення. Прогноз будується за двома останніми прийнятими положеннями $(t_{k-2}, P_{k-2})$, $(t_{k-1}, P_{k-1})$ у припущенні сталої швидкості:
\begin{equation}
\hat{P}_k = P_{k-1} + \mathbf{v} \cdot (t_k - t_{k-1}), \qquad
\mathbf{v} = \frac{P_{k-1} - P_{k-2}}{t_{k-1} - t_{k-2}},
\label{eq:extrapolation}
\end{equation}
де модуль швидкості обмежується значенням $V_{\max}$. За сталого інтервалу між кадрами це збігається з $\hat{P}_k = P_{k-1} + (P_{k-1} - P_{k-2})$.

\textbf{Кінематична перевірка.}\\
Нехай $P_{\text{ост}}$ --- останнє \emph{прийняте} положення, отримане в момент $t_{\text{ост}}$. Нова оцінка $P_k$ приймається, якщо
\begin{equation}
\| P_k - P_{\text{ост}} \|_2 \le D_{\max} = V_{\max} \cdot (t_k - t_{\text{ост}}) + D_{\text{запас}},
\end{equation}
де $V_{\max} = 120$~\px{}/с, $D_{\text{запас}} = 30$~\px{}; за частоти 1~кадр/с $D_{\max} = 150$~\px{}. Порівняння саме з останнім прийнятим положенням означає, що після кількох пропущених кадрів допустимий радіус зростає пропорційно часу. Серед кандидатів, що пройшли фільтр п.~\ref{sec:search}, обирається найкращий за балом, який проходить і кінематичну перевірку. Якщо такого кандидата немає, положення кадру екстраполюється за формулою~(\ref{eq:extrapolation}) (статус \textsc{interpolated}), а відкинута оцінка зберігається для аналізу. Після трьох поспіль кадрів без прийнятого положення трекер переходить у стан \textsc{lost}.

\textbf{Компенсація курсу.}\\
SuperPoint/LightGlue і LoFTR не інваріантні до великих поворотів зображення (кількісне підтвердження наведено в п.~\ref{sec:exp-robust}), а під час розвороту БпЛА кадр повертається відносно карти. Кут <<кадр $\to$ карта>> відомий із гомографій прийнятих кадрів ($\theta$, п.~\ref{sec:search}), тому перед зіставленням кадр повертається на прогнозований кут $\hat{\theta}_k$, і метод зіставлення отримує кадр, орієнтований так само, як карта. Прогноз кута будується аналогічно до~(\ref{eq:extrapolation}) з урахуванням переходу через $\pm 180^\circ$, а кутова швидкість обмежується значенням $45^\circ$/с. Полотно під час повороту розширюється, а відповідності, що потрапили на фон поза кадром, відкидаються. Перетворення <<оригінальний кадр $\to$ робоче зображення>> є композицією масштабування та повороту $T = R(\hat{\theta}_k) \cdot S$, і координати відповідностей переносяться назад точним оберненим перетворенням $T^{-1}$. Тому гомографія, як і раніше, оцінюється в оригінальних пікселях кадру.

Якщо з прогнозованим курсом фрагмент не знайдено, у стані \textsc{tracking} послідовно перевіряються гіпотези $\hat{\theta}_k \pm 45^\circ$ і $\hat{\theta}_k \pm 90^\circ$. У стані \textsc{lost}, де повний перебір карти дорогий, кожен наступний глобальний пошук використовує чергову гіпотезу з циклу $0^\circ, 90^\circ, 180^\circ, -90^\circ$ відносно останнього відомого курсу.

Узагальнено обробка одного кадру в стані \textsc{tracking} складається з таких кроків:
\begin{enumerate}
    \item обчислити прогноз положення $\hat{P}_k$ і курсу $\hat{\theta}_k$;
    \item сформувати пул із не більше ніж 9 фрагментів навколо фрагмента, найближчого до $\hat{P}_k$;
    \item для кожної гіпотези курсу, починаючи з $\hat{\theta}_k$:
    \begin{enumerate}
        \item повернути та масштабувати кадр, обчислити його ознаки;
        \item зіставити кадр із фрагментами пулу, виконати геометричну перевірку, фільтр і ранжування;
        \item якщо найкращий кандидат проходить кінематичну перевірку --- прийняти положення (статус \textsc{tracking}) і завершити обробку кадру;
    \end{enumerate}
    \item якщо жодна гіпотеза не дала прийнятного положення --- екстраполювати положення (статус \textsc{interpolated}); після трьох таких кадрів поспіль перейти в стан \textsc{lost}.
\end{enumerate}

\subsection{Модульна структура та відтворюваність}

Програмну систему організовано як пакет із незалежними модулями (табл.~\ref{tab:modules}) і набір сценаріїв командного рядка для кожного етапу.

\begin{table}[H]
\centering
\caption{Основні модулі програмної системи}
\label{tab:modules}
\small
\begin{tabular}{p{0.36\linewidth} p{0.56\linewidth}}
\toprule
\textbf{Модуль} & \textbf{Призначення} \\
\midrule
\code{preprocessing/tiles.py}, \code{frames.py} & розбиття карти, вилучення кадрів, метадані \\
\code{preprocessing/images.py} & читання/запис, масштабування, поворот, перенесення координат, спотворення для експериментів \\
\code{preprocessing/cache.py} & дисковий кеш ознак фрагментів \\
\code{matching/lightglue.py}, \code{loftr.py} & реалізації методів зі спільним інтерфейсом \\
\code{matching/pipeline.py} & єдиний конвеєр <<кадр + фрагмент $\to$ відповідності $\to$ RANSAC $\to$ метрики>> \\
\code{geometry/homography.py} & RANSAC/MAGSAC++, похибка репроєкції, перевірка правдоподібності \\
\code{evaluation/scoring.py} & фільтр і бал ранжування фрагментів \\
\code{localization/coordinates.py} & перенесення точок на карту, контур огляду \\
\code{localization/search.py}, \code{tracker.py} & пошук фрагмента, трекер послідовності \\
\code{evaluation/metrics.py}, \code{ground\_truth.py} & метрики, еталонні координати \\
\code{visualization/} & рисунки відповідностей, положення та траєкторій \\
\bottomrule
\end{tabular}
\end{table}

Для відтворюваності результатів кожен запуск записує, разом із результатами, файл метаданих з усіма параметрами конвеєра й методу зіставлення, ідентифікатором версії коду та відомостями про обчислювальне середовище. Генератори випадкових чисел RANSAC і спотворень ініціалізуються фіксованими значеннями. Коректність перенесення координат, оцінювання гомографії, критеріїв фільтрації та логіки трекера перевіряється набором із 72 автоматичних тестів на синтетичних даних, для яких точний результат відомий заздалегідь; тести не потребують нейромереж і виконуються за кілька секунд.

\subsection{Візуалізація результатів}

Система формує такі графічні матеріали:
\begin{itemize}
    \item відповідності між кадром і фрагментом карти: геометрично узгоджені --- зеленим, решта --- червоним;
    \item контур огляду камери і центр кадру на ортофотоплані разом із сіткою перевірених фрагментів і обраним фрагментом;
    \item траєкторію польоту на карті з позначенням статусів кадрів і напрямку руху;
    \item графіки похибок, часу обробки та результатів експериментів (розділ~\ref{sec:experiments}).
\end{itemize}
